\documentclass[9pt,a4paper]{article} % Puoi cambiare in report/book se preferisci

% ---------- CODIFICA E LINGUA ----------
\usepackage[T1]{fontenc}
\usepackage[utf8]{inputenc}
\usepackage[italian]{babel}

% ---------- GEOMETRIA PAGINA ----------
\usepackage[a4paper,margin=2.2cm]{geometry}
% Se vuoi ancora più testo per pagina:
% \usepackage[a4paper,margin=2cm]{geometry}

% ---------- FONT E MICROTIPOGRAFIA ----------
\usepackage{lmodern}   % Latin Modern, pulito e adatto alla matematica
\usepackage{microtype} % Migliora la leggibilità (interword spacing, protrusion)

% ---------- MATEMATICA ----------
\usepackage{amsmath,amssymb,amsfonts,amsthm}
\usepackage{mathtools} % Estensione di amsmath
\usepackage{bm}        % Bold math: \bm{...}
\usepackage{physics}   % Comandi utili in fisica (bra-ket, deriv., ecc.)
\usepackage{slashed}   % Notazione slash di Dirac: \slashed{p}
\usepackage{tensor}    % Per indici tensori: \tensor{T}{^\mu_\nu}

\numberwithin{equation}{section} % Equazioni numerate secondo la sezione

% ---------- GRAFICA E COLORI ----------
\usepackage{graphicx}
\usepackage{xcolor}

% ---------- HYPERREF & RIFERIMENTI ----------
\usepackage{hyperref}
\usepackage{cleveref} % Riferimenti intelligenti: \Cref, \cref

\hypersetup{
  colorlinks = true,
  linkcolor  = blue,
  citecolor  = teal,
  urlcolor   = magenta
}

% ---------- FORMATO ARTICOLO SCIENTIFICO ----------
\usepackage{authblk}

% ---------- BIBLIOGRAFIA ----------
% Per citare articoli scientifici con \cite{...} e stampare la bibliografia finale.
\usepackage[
    backend=biber,
    style=numeric,
    sorting=none
]{biblatex}

% Nome del file .bib contenente gli articoli.
\addbibresource{bibliografia.bib}

% ---------- COLONNE MULTIPLE ----------
\usepackage{multicol}

% ---------- AMBIENTI TEOREMI / DEFINIZIONI (OPZIONALE) ----------
\theoremstyle{definition}
\newtheorem{definition}{Definizione}[section]
\theoremstyle{plain}
\newtheorem{theorem}{Teorema}[section]
\newtheorem{proposition}{Proposizione}[section]
\newtheorem{lemma}{Lemma}[section]
\theoremstyle{remark}
\newtheorem{remark}{Osservazione}[section]
\newtheorem{example}{Esempio}[section]

% ---------- INSERIMENTO DI CODICE NEL TESTO ----------
\usepackage{listings}
\usepackage{xcolor}

% Stile generale per i blocchi di codice
\lstdefinestyle{codice}{
    basicstyle=\ttfamily\small,
    numbers=left,
    numberstyle=\tiny,
    stepnumber=1,
    numbersep=8pt,
    frame=single,
    breaklines=true,
    breakatwhitespace=false,
    tabsize=4,
    showstringspaces=false,
    captionpos=b,
    backgroundcolor=\color{gray!8},
    keywordstyle=\color{blue!70!black}\bfseries,
    commentstyle=\color{green!40!black},
    stringstyle=\color{red!60!black}
}

% Imposta questo stile come predefinito
\lstset{style=codice}

% ---------- Per l'indentazione a sinistra dei paragrafi ----------
\setlength{\parindent}{0pt}

% Ecco un esempio di codice inline (chiaramente da inserire all'interno del testo)
% \begin{lstlisting}[language=Python, caption={Simulazione numerica}, label={lst:simulazione}]
% import numpy as np

% N = 100
% x = np.zeros(N)

% for i in range(N):
%     x[i] = i**2
% \end{lstlisting}

% ---------- COMANDI PERSONALIZZATI (ESEMPI) ----------
% Spazio-tempo di Minkowski
\newcommand{\Mink}{\mathbb{M}}
% Operatore d'Alembert
\newcommand{\dal}{\Box}
% Notazione per 4-vettori
\newcommand{\fourvec}[1]{\boldsymbol{#1}}

% ---------- TITOLO DOCUMENTO ----------
\title{Progetto di MAS}

\author[1]{Lorenzo}
\author[1]{Paolo}
\author[1]{Christian}

\affil[1]{Università degli Studi di Torino}

\date{\today}

% \usepackage{xcolor}
% \color{blue}   % oppure red, black, gray, ecc.
% \usepackage{xcolor}
% \definecolor{minvisiblegray}{gray}{0.80} % ~ #C8C8C8
% \color{minvisiblegray}

\begin{document}

% =========================================================
% MODALITÀ ARTICOLO SCIENTIFICO
% =========================================================
% Di base questi comandi sono commentati, quindi il documento resta:
% - senza titolo iniziale;
% - senza abstract;
% - in una sola colonna.

% Per ottenere il formato articolo scientifico:
% 1. scommenta \maketitle;
% 2. scommenta l'ambiente abstract;
% 3. scommenta \twocolumn prima del testo principale;
% 4. scommenta \printbibliography alla fine.

% \maketitle

% \begin{abstract}
% In questo lavoro si presenta ...
% Qui va inserito un breve riassunto dell'articolo: contesto, obiettivo,
% metodo utilizzato e risultati principali.
% \end{abstract}

% \vspace{0.5cm}

\subsection{Estensioni, specie di agenti e variabili principali del modello}

\begin{lstlisting}[caption={Dichiarazione delle estensioni, delle specie e delle variabili principali}]
extensions [ fetch table csv ]

breed [ agents-a agent-a ]
breed [ agents-b agent-b ]
breed [ agents-c agent-c ]

globals [
  prices
  day-index

  ;; DATI REALI (STARTER PER LA CORRENTE)
  btc-price
  btc-initial-price
  btc-ath-price
  btc-prev-price
  btc-return
  btc-volatility

  ;; PREZZO E VOLATILITA' STIMATI (ENDOGENI)
  model-price
  model-return
  model-volatility
  model-trading-volume
  real-trading-volume-usd

  ;; PARAMETRI CORRENTE STARTER E SWITCH DINAMICO
  current-disabled?
  ;; pearson-threshold viene gestito dallo slider nell'interfaccia
  current-hype                  ;; Livello di stress/hype di mercato

  ;; SCALING E ISTERESI VOLUMI
  real-volume-initial
  real-volume-ath
  volume-hysteresis-state

  ;; STORICO E CORRELAZIONE VOLUMI (PEARSON)
  model-volume-history
  real-volume-history
  volume-correlation
  volume-correlation-history
  mean-volume-correlation
  above-threshold-counter       ;; Conta i tick consecutivi sopra pearson-threshold

  ;; STORICO E CORRELAZIONE PREZZI (PEARSON)
  model-price-history
  real-price-history
  price-correlation
  price-correlation-history
  mean-price-correlation

  op-history
  ya-history
  yb-history
  lambda-instability
  heterogeneity-index-value
  phase-correlation-ab
]

turtles-own [
  flockmates
  bitcoin-held
  cash-balance
  velocity-x
  velocity-y
]
\end{lstlisting}

Questa prima parte del codice definisce la struttura generale del modello. In
particolare, vengono dichiarate le estensioni utilizzate, le specie di agenti e
le variabili globali o individuali che saranno impiegate nelle procedure
successive.

La riga

\[
\texttt{extensions [ fetch table csv ]}
\]

carica tre estensioni di NetLogo. L'estensione \texttt{fetch} permette di
effettuare richieste verso risorse esterne, ad esempio scaricando dati da una
API. In questo modello essa è utile per ottenere dati reali relativi a Bitcoin.
L'estensione \texttt{table} consente di manipolare strutture dati del tipo
chiave-valore, particolarmente utili quando si ricevono dati in formato
strutturato. Infine, l'estensione \texttt{csv} permette di leggere o scrivere
dati in formato CSV, cioè in un formato tabellare comunemente usato per salvare
serie storiche o risultati di simulazione.

Le tre righe successive definiscono tre specie di agenti tramite il costrutto
\texttt{breed}:

\[
\texttt{breed [ agents-a agent-a ]},
\qquad
\texttt{breed [ agents-b agent-b ]},
\qquad
\texttt{breed [ agents-c agent-c ]}.
\]

In NetLogo, una \texttt{breed} rappresenta una sottopopolazione di turtles. Il
primo nome all'interno della parentesi quadra è il nome plurale della specie,
mentre il secondo è il nome singolare. Ad esempio, \texttt{agents-a} indica
l'insieme degli agenti della prima specie, mentre \texttt{agent-a} indica un
singolo agente di quella specie.

L'introduzione di più specie permette di distinguere gruppi differenti di agenti.
Dal punto di vista modellistico, tali gruppi possono rappresentare categorie
diverse di operatori, strategie differenti, popolazioni eterogenee oppure classi
di agenti con comportamenti o parametri differenti. Anche se appartengono a
specie diverse, tutti questi agenti rimangono comunque turtles: un comando come

\[
\texttt{ask turtles [ ... ]}
\]

agisce quindi su tutte le specie contemporaneamente.

Il blocco \texttt{globals} contiene le variabili globali del modello. Una
variabile globale non appartiene a un singolo agente, ma descrive una proprietà
dell'intero sistema o una grandezza accessibile da tutte le procedure.

Le prime due variabili sono

\[
\texttt{prices},
\qquad
\texttt{day-index}.
\]

La variabile \texttt{prices} contiene la serie storica dei prezzi scaricati o
letti dal modello. La variabile \texttt{day-index} indica invece quale elemento
della serie storica viene utilizzato al passo corrente della simulazione. Essa
funziona quindi come indice temporale: aumentando \texttt{day-index}, il modello
procede lungo la sequenza dei dati di mercato.

La sezione successiva riguarda i dati reali di Bitcoin:

\[
\texttt{btc-price},
\quad
\texttt{btc-initial-price},
\quad
\texttt{btc-ath-price},
\quad
\texttt{btc-prev-price},
\quad
\texttt{btc-return},
\quad
\texttt{btc-volatility}.
\]

Bitcoin è una criptovaluta, cioè un'attività digitale scambiata su mercati
elettronici. Come per un'azione o per un altro strumento finanziario, il suo
prezzo varia nel tempo in base all'incontro tra domanda e offerta.

La variabile \texttt{btc-price} rappresenta il prezzo reale corrente di Bitcoin
nel dato temporale considerato. La variabile \texttt{btc-initial-price} memorizza
il prezzo iniziale della serie utilizzata, mentre \texttt{btc-ath-price}
rappresenta il massimo storico raggiunto nella finestra di dati considerata. La
sigla ATH significa \textit{All-Time High}, cioè massimo storico o massimo
raggiunto da una grandezza in un certo intervallo di osservazione.

La variabile \texttt{btc-prev-price} contiene il prezzo precedente di Bitcoin ed
è necessaria per calcolare il rendimento. In finanza, il rendimento misura la
variazione relativa del prezzo tra due istanti consecutivi. Se \(P(t)\) indica
il prezzo corrente e \(P(t-1)\) il prezzo precedente, il rendimento può essere
definito come

\[
r(t)=\frac{P(t)-P(t-1)}{P(t-1)}.
\]

Nel modello, il rendimento reale di Bitcoin viene memorizzato nella variabile

\[
\texttt{btc-return}.
\]

La variabile \texttt{btc-volatility} rappresenta invece la volatilità reale di
Bitcoin. In ambito finanziario, la volatilità misura l'intensità delle variazioni
di prezzo. Un mercato molto volatile è un mercato in cui il prezzo cambia molto
rapidamente o con variazioni ampie. In una versione semplificata, la volatilità
può essere calcolata come valore assoluto del rendimento:

\[
V(t)=|r(t)|.
\]

Questa scelta elimina il segno della variazione: una forte crescita e una forte
diminuzione del prezzo producono entrambe un valore elevato di volatilità.

Il blocco successivo introduce le variabili endogene del modello:

\[
\texttt{model-price},
\quad
\texttt{model-return},
\quad
\texttt{model-volatility},
\quad
\texttt{model-trading-volume},
\quad
\texttt{real-trading-volume-usd}.
\]

Una variabile si dice endogena quando viene prodotta internamente dal modello. In
questo caso, \texttt{model-price} rappresenta il prezzo generato dalla dinamica
simulata degli agenti, mentre \texttt{btc-price} rappresenta un prezzo reale
esterno. La distinzione tra prezzo reale e prezzo del modello è importante:
il primo deriva dai dati empirici, il secondo deriva invece dall'interazione tra
gli agenti.

La variabile \texttt{model-return} rappresenta il rendimento del prezzo endogeno,
calcolato confrontando il prezzo prodotto dal modello con il suo valore
precedente. Analogamente, \texttt{model-volatility} misura l'intensità delle
variazioni del prezzo simulato.

La variabile \texttt{model-trading-volume} rappresenta il volume di scambio
generato dagli agenti del modello. In finanza, il volume di scambio indica la
quantità complessiva di un'attività finanziaria comprata o venduta in un certo
intervallo di tempo. Nel caso del modello, il volume non deriva dal mercato reale,
ma dal numero di operazioni effettuate dagli agenti.

La variabile \texttt{real-trading-volume-usd} rappresenta invece il volume reale
del mercato, espresso in dollari. Il suffisso \texttt{usd} indica che il valore è
misurato in dollari statunitensi. Questa variabile permette di confrontare il
volume simulato con il volume reale osservato.

La sezione dedicata alla corrente e allo switch dinamico contiene le variabili

\[
\texttt{current-disabled?},
\qquad
\texttt{current-hype}.
\]

La variabile \texttt{current-disabled?} è una variabile booleana, cioè una
variabile che può assumere valore vero o falso. Il punto interrogativo finale è
una convenzione frequente in NetLogo per indicare variabili o reporter logici.
Essa serve a stabilire se la corrente esterna deve essere disattivata.

La corrente rappresenta una spinta o un campo esterno che influenza il movimento
degli agenti. Nelle prime fasi del modello, questa corrente può essere guidata
dai dati reali del mercato, ad esempio dalla volatilità di Bitcoin. In seguito,
il modello può decidere di disattivare o modificare questa dipendenza esterna in
base alla qualità dell'accordo tra il mercato reale e il mercato simulato.

La variabile \texttt{current-hype} rappresenta il livello di stress o hype di
mercato. Dal punto di vista economico, il termine \textit{hype} indica una fase
di forte attenzione, entusiasmo o agitazione collettiva intorno a un'attività
finanziaria. In un mercato con alto hype, gli agenti possono essere più reattivi,
più propensi a operare e più sensibili alle variazioni di prezzo o volume.

Il commento relativo a \texttt{pearson-threshold} indica che tale soglia non è
dichiarata come variabile globale nel codice, perché viene gestita tramite uno
slider dell'interfaccia. Essa rappresenta una soglia utilizzata per stabilire se
la correlazione tra dati reali e dati simulati è sufficientemente elevata.

La sezione successiva riguarda lo scaling e l'isteresi dei volumi:

\[
\texttt{real-volume-initial},
\qquad
\texttt{real-volume-ath},
\qquad
\texttt{volume-hysteresis-state}.
\]

La variabile \texttt{real-volume-initial} memorizza il volume reale iniziale,
mentre \texttt{real-volume-ath} memorizza il massimo volume reale osservato nella
finestra di dati considerata. Anche qui, \texttt{ath} indica un massimo storico o
un massimo relativo all'intervallo analizzato.

Il concetto di scaling indica una trasformazione usata per portare grandezze di
ordini di grandezza diversi su scale confrontabili. Questo è utile perché il
volume reale di mercato può essere enormemente più grande del volume prodotto da
un numero limitato di agenti simulati.

La variabile \texttt{volume-hysteresis-state} rappresenta lo stato di isteresi
dei volumi. In generale, si parla di isteresi quando lo stato attuale di un
sistema non dipende solo dal valore corrente di una variabile, ma anche dalla sua
storia precedente. Dal punto di vista economico, ciò può rappresentare un effetto
di memoria del mercato: un periodo di forte volume o forte stress può continuare
a influenzare il comportamento degli agenti anche dopo che il segnale originario
si è ridotto.

Le variabili successive memorizzano lo storico dei volumi e la loro correlazione:

\[
\texttt{model-volume-history},
\quad
\texttt{real-volume-history},
\quad
\texttt{volume-correlation},
\quad
\texttt{volume-correlation-history},
\quad
\texttt{mean-volume-correlation},
\quad
\texttt{above-threshold-counter}.
\]

Le variabili \texttt{model-volume-history} e \texttt{real-volume-history}
contengono rispettivamente la sequenza temporale dei volumi prodotti dal modello
e dei volumi reali. Questi storici sono necessari per confrontare la dinamica
simulata con quella empirica.

La variabile \texttt{volume-correlation} contiene il coefficiente di correlazione
di Pearson tra i volumi del modello e i volumi reali. La correlazione di Pearson
misura il grado di relazione lineare tra due serie di dati. Se \(X\) e \(Y\)
sono due serie temporali, il coefficiente di Pearson può essere scritto come

\[
\rho_{X,Y}
=
\frac{
\sum_{i=1}^{N}(X_i-\bar{X})(Y_i-\bar{Y})
}{
\sqrt{\sum_{i=1}^{N}(X_i-\bar{X})^2}
\sqrt{\sum_{i=1}^{N}(Y_i-\bar{Y})^2}
}.
\]

Il valore di \(\rho\) appartiene all'intervallo \([-1,1]\). Se \(\rho\) è vicino
a \(1\), le due serie tendono a muoversi nella stessa direzione. Se \(\rho\) è
vicino a \(-1\), tendono a muoversi in direzioni opposte. Se \(\rho\) è vicino
a \(0\), non emerge una relazione lineare significativa.

La variabile \texttt{volume-correlation-history} conserva lo storico dei valori
di correlazione calcolati nel tempo, mentre \texttt{mean-volume-correlation}
memorizza una correlazione media. Infine,
\texttt{above-threshold-counter} conta per quanti tick consecutivi la
correlazione rimane sopra la soglia \texttt{pearson-threshold}. Questa variabile
è utile per evitare decisioni basate su oscillazioni momentanee: il modello può
richiedere che la correlazione rimanga alta per più passi temporali prima di
modificare il proprio regime dinamico.

Una struttura analoga viene usata per i prezzi:

\[
\texttt{model-price-history},
\quad
\texttt{real-price-history},
\quad
\texttt{price-correlation},
\quad
\texttt{price-correlation-history},
\quad
\texttt{mean-price-correlation}.
\]

In questo caso il confronto non riguarda i volumi, ma i prezzi. La variabile
\texttt{model-price-history} conserva la serie storica dei prezzi endogeni
prodotti dal modello, mentre \texttt{real-price-history} conserva la serie dei
prezzi reali di Bitcoin. La variabile \texttt{price-correlation} misura la
correlazione di Pearson tra queste due serie.

Il confronto tra prezzo simulato e prezzo reale permette di valutare se la
dinamica interna del modello riesce a riprodurre, almeno in parte, l'andamento
osservato nel mercato reale. Anche in questo caso è importante ricordare che una
correlazione elevata non implica causalità: indica soltanto che le due serie
tendono a muoversi in modo linearmente simile.

Le ultime variabili globali sono

\[
\texttt{op-history},
\quad
\texttt{ya-history},
\quad
\texttt{yb-history},
\quad
\texttt{lambda-instability},
\quad
\texttt{heterogeneity-index-value},
\quad
\texttt{phase-correlation-ab}.
\]

La variabile \texttt{op-history} può essere interpretata come lo storico del
parametro d'ordine. Il parametro d'ordine misura il grado di allineamento
collettivo degli agenti: valori vicini a \(1\) indicano un sistema molto
allineato, mentre valori vicini a \(0\) indicano un sistema disordinato.

Le variabili \texttt{ya-history} e \texttt{yb-history} memorizzano storici legati
alle coordinate verticali o alle dinamiche medie di specifiche specie di agenti.
Poiché nel modello la coordinata \(y\) può essere associata alla propensione ad
acquistare o vendere, questi storici possono essere utilizzati per studiare
l'evoluzione comportamentale dei gruppi.

La variabile \texttt{lambda-instability} rappresenta una misura di instabilità
del sistema. Il simbolo \(\lambda\) viene spesso usato per indicare un parametro
che misura l'intensità di una crescita, di una divergenza o di una sensibilità
dinamica. Nel contesto del modello, questa variabile può essere usata per
quantificare quanto il sistema sia stabile o instabile nel tempo.

La variabile \texttt{heterogeneity-index-value} rappresenta un indice di
eterogeneità. L'eterogeneità misura quanto gli agenti o i gruppi differiscono tra
loro. Un valore elevato può indicare che specie diverse assumono comportamenti
molto differenti, mentre un valore basso può indicare una maggiore omogeneità
della popolazione.

Infine, \texttt{phase-correlation-ab} misura una correlazione di fase tra le
specie \texttt{agents-a} e \texttt{agents-b}. In termini qualitativi, questa
grandezza può indicare se due gruppi oscillano o si muovono in modo coordinato,
oppure se seguono dinamiche sfasate.

Il blocco \texttt{turtles-own} dichiara invece le variabili individuali degli
agenti:

\[
\texttt{flockmates},
\quad
\texttt{bitcoin-held},
\quad
\texttt{cash-balance},
\quad
\texttt{velocity-x},
\quad
\texttt{velocity-y}.
\]

A differenza delle variabili globali, queste variabili appartengono
separatamente a ciascun agente. Ogni turtle possiede quindi un proprio valore di
\texttt{bitcoin-held}, un proprio valore di \texttt{cash-balance} e così via.

La variabile \texttt{flockmates} contiene l'insieme degli agenti vicini con cui
la turtle interagisce nella regola di flocking. Essa è fondamentale per
l'allineamento locale, perché permette a ogni agente di individuare i vicini
rispetto ai quali aggiornare la propria direzione.

La variabile \texttt{bitcoin-held} indica quanti Bitcoin sono posseduti da un
agente. Essa rappresenta quindi la posizione dell'agente sull'attività
finanziaria. La variabile \texttt{cash-balance} rappresenta invece la liquidità
disponibile, cioè la quantità di denaro che l'agente può usare per acquistare
Bitcoin

\subsection{Helper per la gestione sicura delle penne dei plot}

\begin{lstlisting}[caption={Procedura di controllo e creazione delle penne dei plot}]
;;;;;;;;;;;;;;;;;;;;;;;;;;;;;;;;;;;;;;
;; HELPER PER PLOT SICURI
;;;;;;;;;;;;;;;;;;;;;;;;;;;;;;;;;;;;;;

to assicura-penna [ nome-penna colore ]
  carefully [
    set-current-plot-pen nome-penna
  ] [
    create-temporary-plot-pen nome-penna
    set-plot-pen-color colore
  ]
end
\end{lstlisting}

La procedura \texttt{assicura-penna} è una procedura di supporto pensata per
rendere più robusta la gestione dei grafici. Essa riceve due argomenti:

\[
\texttt{nome-penna},
\qquad
\texttt{colore}.
\]

Il primo indica il nome della penna del plot che si vuole utilizzare, mentre il
secondo indica il colore da assegnare alla penna nel caso in cui questa debba
essere creata.

Il blocco principale è costruito tramite il comando

\[
\texttt{carefully [ ... ] [ ... ]}.
\]

In NetLogo, \texttt{carefully} permette di provare a eseguire un blocco di
codice e, se questo genera un errore, eseguire un secondo blocco alternativo.
In questo caso, il modello prova prima a selezionare una penna già esistente:

\[
\texttt{set-current-plot-pen nome-penna}.
\]

Se la penna esiste, viene semplicemente selezionata. Se invece non esiste,
NetLogo genererebbe normalmente un errore. Grazie a \texttt{carefully}, tale
errore viene intercettato e viene eseguito il blocco alternativo:

\[
\texttt{create-temporary-plot-pen nome-penna}
\]

seguito da

\[
\texttt{set-plot-pen-color colore}.
\]

In questo modo, se la penna richiesta non è già presente nel plot, viene creata
automaticamente e le viene assegnato il colore specificato.

Questa procedura è utile quando il modello aggiorna molti grafici o molte curve
diverse. Invece di dover creare manualmente tutte le penne dall'interfaccia di
NetLogo, il codice controlla automaticamente se una penna esiste e, se manca, la
crea. Questo rende i plot più flessibili e riduce il rischio di errori dovuti a
nomi di penne non ancora definiti.

\subsection{Inizializzazione del portafoglio degli agenti}

\begin{lstlisting}[caption={Procedura di inizializzazione del portafoglio individuale}]
;;;;;;;;;;;;;;;;;;;;;;;;;;;;;;;;;;;;;;
;; INIZIALIZZAZIONE PORTAFOGLIO
;;;;;;;;;;;;;;;;;;;;;;;;;;;;;;;;;;;;;;

to inizializza-portafoglio
  set bitcoin-held random 6
  set cash-balance 50000 + random 1450001
end
\end{lstlisting}

La procedura \texttt{inizializza-portafoglio} assegna a ogni agente una dotazione
finanziaria iniziale. Essa viene normalmente richiamata durante la creazione
degli agenti nella procedura \texttt{setup}.

La prima istruzione è

\[
\texttt{set bitcoin-held random 6}.
\]

La variabile \texttt{bitcoin-held} rappresenta il numero di Bitcoin posseduti
dall'agente. Il comando \texttt{random 6} genera un numero intero casuale tra
\(0\) e \(5\). Di conseguenza, ogni agente inizia la simulazione con una quantità
casuale di Bitcoin compresa nell'intervallo

\[
[0,5].
\]

La seconda istruzione è

\[
\texttt{set cash-balance 50000 + random 1450001}.
\]

La variabile \texttt{cash-balance} rappresenta la liquidità disponibile
dell'agente, cioè la quantità di denaro che può essere usata per effettuare
acquisti. Il termine \texttt{random 1450001} genera un numero intero casuale tra
\(0\) e \(1\,450\,000\). Sommando \(50\,000\), si ottiene una liquidità iniziale
compresa tra

\[
50\,000
\quad \text{e} \quad
1\,500\,000.
\]

Questa inizializzazione introduce eterogeneità economica tra gli agenti: alcuni
partono con molti Bitcoin e molta liquidità, altri con meno disponibilità.

Dal punto di vista economico, il portafoglio di un agente è composto da due
elementi: una parte detenuta in Bitcoin e una parte detenuta in liquidità. La
liquidità indica la capacità immediata di acquistare nuovi Bitcoin, mentre
\texttt{bitcoin-held} indica la quantità dell'attività finanziaria già posseduta.
Questa distinzione è necessaria per impedire operazioni non realistiche, come
comprare senza denaro disponibile o vendere Bitcoin non posseduti.

\subsection{Posizionamento iniziale degli agenti in una regione centrale}

\begin{lstlisting}[caption={Procedura di spawn circolare centrale}]
;;;;;;;;;;;;;;;;;;;;;;;;;;;;;;;;;;;;;;
;; SPAWN CIRCOLARE CENTRALE
;;;;;;;;;;;;;;;;;;;;;;;;;;;;;;;;;;;;;;

to posiziona-al-centro
  let angolo random-float 360
  let raggio random-float spawn-radius

  setxy (raggio * cos angolo)
        (raggio * sin angolo)
end
\end{lstlisting}

La procedura \texttt{posiziona-al-centro} assegna a ogni agente una posizione
iniziale casuale all'interno di una regione circolare centrata nell'origine del
mondo NetLogo. Essa viene utilizzata al momento della creazione degli agenti, in
modo che la popolazione non venga distribuita casualmente in tutto lo spazio, ma
nasca inizialmente in una zona centrale controllata.

La variabile locale

\[
\texttt{angolo}
\]

viene definita tramite

\[
\texttt{let angolo random-float 360}.
\]

Il comando \texttt{random-float 360} genera un numero reale casuale compreso tra
\(0\) e \(360\). Questo valore rappresenta l'angolo polare della posizione
iniziale dell'agente. In altre parole, stabilisce la direzione, rispetto
all'origine, lungo la quale verrà collocato l'agente.

La variabile locale

\[
\texttt{raggio}
\]

viene invece definita tramite

\[
\texttt{let raggio random-float spawn-radius}.
\]

Il parametro \texttt{spawn-radius} controlla la dimensione della regione iniziale
di spawn. Il comando \texttt{random-float spawn-radius} genera un raggio casuale
compreso tra \(0\) e \texttt{spawn-radius}. Quindi gli agenti vengono collocati
entro un disco di raggio massimo \texttt{spawn-radius}.

La posizione cartesiana viene poi ottenuta a partire dalle coordinate polari:

\[
x = r \cos(\theta),
\qquad
y = r \sin(\theta).
\]

Nel codice NetLogo questa trasformazione è implementata da

\[
\texttt{setxy (raggio * cos angolo) (raggio * sin angolo)}.
\]

Il comando \texttt{setxy} assegna simultaneamente la coordinata \(x\) e la
coordinata \(y\) dell'agente. Poiché il centro del disco è l'origine, non è
necessario aggiungere alcuna traslazione alle coordinate.

Dal punto di vista modellistico, questa procedura permette di studiare
l'evoluzione di un gruppo inizialmente concentrato. Gli agenti partono da una
regione comune e successivamente si disperdono, si allineano o si separano in
base alle regole di flocking, repulsione, corrente e trading definite nel
modello.

\subsection{Procedura di inizializzazione del modello: \texttt{setup}}

\begin{lstlisting}[caption={Procedura di setup del modello}]
;;;;;;;;;;;;;;;;;;;;;;;;;;;;;;;;;;;;;;
;; SETUP
;;;;;;;;;;;;;;;;;;;;;;;;;;;;;;;;;;;;;;

to setup

  clear-all

  carefully [
    import-drawing "grafico_buy_sell.png"
  ] [ ]

  ;; GARANZIA SOGLIA PEARSON VALIDA DA SLIDER
  if not is-number? pearson-threshold [
    set pearson-threshold 0.5
  ]

  set current-disabled? false    ;; Inizia sempre in Fase 1
  set current-hype 0

  if is-number? price-sensitivity = false [
    set price-sensitivity 0.01
  ]

  set day-index 0
  set btc-prev-price 0
  set btc-price 0
  set btc-initial-price 0
  set btc-ath-price 0
  set btc-return 0
  set btc-volatility 0

  set model-price 0
  set model-return 0
  set model-volatility 0

  set model-trading-volume 0
  set real-trading-volume-usd 0
  set real-volume-initial 0
  set real-volume-ath 0
  set volume-hysteresis-state 0

  ;; INIZIALIZZAZIONE CORRELAZIONE VOLUMI
  set volume-correlation 0
  set volume-correlation-history []
  set mean-volume-correlation 0
  set above-threshold-counter 0

  ;; INIZIALIZZAZIONE CORRELAZIONE PREZZI
  set model-price-history []
  set real-price-history []
  set price-correlation 0
  set price-correlation-history []
  set mean-price-correlation 0

  set prices []
  set model-volume-history []
  set real-volume-history []

  set op-history []
  set ya-history []
  set yb-history []

  set lambda-instability 0
  set heterogeneity-index-value 0
  set phase-correlation-ab 0

  if trade-interval <= 0 [
    user-message "Attenzione: trade-interval deve essere > 0. Impostato a 1."
    set trade-interval 1
  ]

  scarica-dati

  ;; CREAZIONE AGENTI
  create-agents-a population-a [
    set color yellow
    set size 10
    posiziona-al-centro
    set heading random 360
    set velocity-x cos heading
    set velocity-y sin heading
    set flockmates no-turtles
    inizializza-portafoglio
  ]

  create-agents-b population-b [
    set color red
    set size 10
    posiziona-al-centro
    set heading random 360
    set velocity-x cos heading
    set velocity-y sin heading
    set flockmates no-turtles
    inizializza-portafoglio
  ]

  create-agents-c population-c [
    set color cyan
    set size 10
    posiziona-al-centro
    set heading random 360
    set velocity-x cos heading
    set velocity-y sin heading
    set flockmates no-turtles
    inizializza-portafoglio
  ]

  ;; RETE FISSA
  ask agents-a [
    let n min (list number-of-followers (count other agents-a))
    if n > 0 [ set flockmates n-of n other agents-a ]
  ]

  ask agents-b [
    let n min (list number-of-followers (count other agents-b))
    if n > 0 [ set flockmates n-of n other agents-b ]
  ]

  ask agents-c [
    let n min (list number-of-followers (count other agents-c))
    if n > 0 [ set flockmates n-of n other agents-c ]
  ]

  ;; INIZIALIZZAZIONE PRIMO GIORNO REALE E MODELLO
  if is-list? prices and length prices > 0 [
    let record item 0 prices
    set btc-price item 1 record
    set btc-prev-price btc-price
    set btc-initial-price btc-price
    set btc-ath-price btc-price
    set day-index 1
  ]

  set model-price btc-price
  set model-return 0
  set model-volatility 0

  print (word "=== SETUP COMPLETATO ===")
  print (word "FASE 1 ATTIVA - Switch automatico a Fase 2 quando r_medio(volumi) >= " pearson-threshold " per piu' di 5 tick consecutivi.")

  reset-ticks

end
\end{lstlisting}

La procedura \texttt{setup} inizializza l'intero modello e prepara sia la parte
finanziaria sia la parte agent-based della simulazione. Il comando
\texttt{clear-all} elimina lo stato precedente del modello, cancellando agenti,
plot, variabili temporanee e configurazioni residue. In questo modo ogni nuova
esecuzione parte da una condizione pulita.

Subito dopo, il codice prova a importare l'immagine
\texttt{grafico\_buy\_sell.png} tramite il blocco

\[
\texttt{carefully [ import-drawing ... ] [ ]}
\]

Il comando \texttt{carefully} permette di evitare che il modello si blocchi nel
caso in cui il file immagine non sia presente. Se l'importazione fallisce, il
secondo blocco vuoto viene eseguito e la simulazione continua comunque.

La parte successiva controlla che alcuni parametri dell'interfaccia abbiano
valori numerici validi. In particolare, \texttt{pearson-threshold} viene impostato
a \(0.5\) se non è un numero. Questa variabile rappresenta la soglia di
correlazione di Pearson richiesta per considerare sufficientemente simili due
serie temporali, ad esempio i volumi reali e i volumi generati dal modello.
Analogamente, \texttt{price-sensitivity} viene impostata a \(0.01\) se non è
definita correttamente. Tale parametro controlla quanto il prezzo endogeno del
modello reagisce agli squilibri generati dagli agenti.

La variabile \texttt{current-disabled?} viene inizializzata a \texttt{false}.
Questo indica che, all'inizio, la corrente esterna è attiva. Il commento specifica
che il modello parte sempre dalla \emph{Fase 1}, cioè da una fase in cui il
sistema è ancora guidato, almeno in parte, dai dati reali. La variabile
\texttt{current-hype} viene posta a zero e rappresenta il livello iniziale di
hype o stress di mercato.

Segue l'inizializzazione delle variabili relative al mercato reale di Bitcoin:
\texttt{btc-price}, \texttt{btc-prev-price}, \texttt{btc-initial-price},
\texttt{btc-ath-price}, \texttt{btc-return} e \texttt{btc-volatility}. Il prezzo
di Bitcoin rappresenta il valore dell'attività finanziaria nel tempo; il
rendimento misura la variazione relativa del prezzo tra due istanti consecutivi;
la volatilità misura invece l'intensità delle variazioni di prezzo. Il valore
\texttt{btc-ath-price} memorizza il massimo prezzo raggiunto nella finestra di
osservazione considerata.

Vengono poi inizializzate le variabili endogene del modello:
\texttt{model-price}, \texttt{model-return} e \texttt{model-volatility}. A
differenza delle variabili precedenti, queste non provengono direttamente dai
dati reali, ma sono generate dalla dinamica interna della simulazione. Il prezzo
endogeno è quindi il prezzo prodotto dal comportamento degli agenti.

La procedura azzera anche le variabili relative ai volumi:
\texttt{model-trading-volume}, \texttt{real-trading-volume-usd},
\texttt{real-volume-initial}, \texttt{real-volume-ath} e
\texttt{volume-hysteresis-state}. In economia, il volume di scambio misura la
quantità o il controvalore di un'attività comprata e venduta in un certo
intervallo di tempo. Nel modello si distinguono i volumi reali, provenienti dal
mercato, e i volumi endogeni, generati dagli scambi degli agenti.

Le liste storiche e le variabili di correlazione vengono inizializzate come liste
vuote o valori nulli. In particolare, \texttt{model-volume-history} e
\texttt{real-volume-history} conterranno gli storici dei volumi simulati e reali,
mentre \texttt{model-price-history} e \texttt{real-price-history} conterranno gli
storici dei prezzi. Le variabili \texttt{volume-correlation} e
\texttt{price-correlation} misureranno la correlazione di Pearson tra le serie
simulate e quelle reali. Il contatore \texttt{above-threshold-counter} serve a
verificare per quanti tick consecutivi la correlazione dei volumi rimane sopra
la soglia fissata.

La procedura inizializza anche alcune grandezze dinamiche legate al comportamento
collettivo degli agenti: \texttt{op-history}, \texttt{ya-history},
\texttt{yb-history}, \texttt{lambda-instability},
\texttt{heterogeneity-index-value} e \texttt{phase-correlation-ab}. Queste
variabili permettono di studiare il grado di ordine, l'instabilità, l'eterogeneità
tra gruppi e la relazione di fase tra specie diverse.

Prima di scaricare i dati, il codice controlla che \texttt{trade-interval} sia
positivo. Questo parametro stabilisce ogni quante iterazioni avviene il trading.
Se il valore è minore o uguale a zero, NetLogo mostra un messaggio di avviso e
imposta automaticamente \texttt{trade-interval} a \(1\), evitando divisioni o
condizioni temporali non valide.

Il comando

\[
\texttt{scarica-dati}
\]

richiama la procedura che carica la serie dei prezzi reali di Bitcoin. Dopo il
download, il modello crea le tre popolazioni di agenti:
\texttt{agents-a}, \texttt{agents-b} e \texttt{agents-c}. Le tre specie vengono
distinte tramite il colore: giallo, rosso e ciano. Ogni agente viene posizionato
nella regione centrale tramite \texttt{posiziona-al-centro}, riceve una direzione
iniziale casuale e le corrispondenti componenti di velocità:

\[
v_x = \cos(\theta),
\qquad
v_y = \sin(\theta),
\]

dove \(\theta\) è il valore di \texttt{heading}. Ogni agente viene inoltre
inizializzato senza vicini di flocking e riceve un portafoglio iniziale tramite
\texttt{inizializza-portafoglio}.

La sezione denominata \emph{rete fissa} assegna a ciascun agente un insieme
iniziale di \texttt{flockmates} appartenenti alla stessa specie. Per ogni gruppo,
il numero di vicini scelti è il minimo tra \texttt{number-of-followers} e il
numero effettivo di altri agenti disponibili nella stessa specie. Il comando
\texttt{n-of} seleziona casualmente \(n\) agenti. In questo modo il flocking non
dipende necessariamente da una vicinanza spaziale istantanea, ma da una rete di
relazioni interne alla specie.

Infine, se la lista \texttt{prices} contiene dati validi, il primo record viene
usato per inizializzare il primo prezzo reale di Bitcoin. Il prezzo corrente,
il prezzo precedente, il prezzo iniziale e il massimo iniziale vengono tutti
posti uguali al primo prezzo disponibile. L'indice temporale \texttt{day-index}
viene poi impostato a \(1\), così che l'aggiornamento successivo legga il record
seguente.

Il prezzo endogeno \texttt{model-price} viene inizializzato uguale al prezzo
reale iniziale. Questa scelta permette al modello simulato di partire dallo stesso
livello del mercato reale, rendendo più significativo il confronto successivo tra
prezzo reale e prezzo generato dagli agenti.

La procedura termina stampando alcuni messaggi informativi nella console e
richiamando \texttt{reset-ticks}, che azzera il tempo della simulazione.

\subsection{Hype di mercato e stubbornness dinamica}

\begin{lstlisting}[caption={Calcolo dell'hype e della stubbornness dinamica}]
;;;;;;;;;;;;;;;;;;;;;;;;;;;;;;;;;;;;;;
;; HYPE E STUBBORNNESS DINAMICA
;;;;;;;;;;;;;;;;;;;;;;;;;;;;;;;;;;;;;;

to-report calcola-hype
  report (abs model-return) + model-volatility
end

to-report dynamic-stubbornness
  let base-s 0.95
  let min-s 0.05
  let raw-s base-s - (hype-sensitivity * current-hype)
  report max (list min-s (min (list base-s raw-s)))
end
\end{lstlisting}

Questo blocco introduce due reporter legati alla reattività degli agenti rispetto
alle condizioni del mercato simulato. Il primo reporter, \texttt{calcola-hype},
definisce una misura sintetica dell'hype o stress di mercato:

\[
H(t)=|r_m(t)|+\sigma_m(t),
\]

dove \(r_m(t)\) è il rendimento del prezzo endogeno del modello e
\(\sigma_m(t)\) è la volatilità endogena. Il termine \(|r_m(t)|\) misura
l'intensità della variazione del prezzo, indipendentemente dal segno, mentre
\texttt{model-volatility} rappresenta il grado di instabilità del prezzo
simulato.

Dal punto di vista economico, l'hype può essere interpretato come una misura
aggregata della tensione del mercato. Quando il prezzo cambia molto e la
volatilità è elevata, il mercato viene considerato più instabile, più reattivo o
più emotivamente carico. In questo senso, \texttt{current-hype} rappresenta una
variabile comportamentale che traduce l'instabilità del prezzo in una maggiore
sensibilità degli agenti.

Il secondo reporter, \texttt{dynamic-stubbornness}, calcola una stubbornness
dinamica, cioè una misura della tendenza degli agenti a conservare il proprio
orientamento o comportamento precedente. Il valore di partenza è

\[
s_{\mathrm{base}}=0.95,
\]

mentre il valore minimo ammesso è

\[
s_{\min}=0.05.
\]

La quantità intermedia

\[
s_{\mathrm{raw}}
=
s_{\mathrm{base}}
-
\texttt{hype-sensitivity}\cdot \texttt{current-hype}
\]

riduce la stubbornness quando aumenta l'hype di mercato. Quindi, in condizioni
di mercato calme, gli agenti rimangono più stabili e meno reattivi; in condizioni
di mercato stressato, diventano più sensibili e più disposti a modificare il loro
stato.

L'ultima riga,

\[
\texttt{report max (list min-s (min (list base-s raw-s)))},
\]

impone che la stubbornness resti sempre compresa tra \texttt{min-s} e
\texttt{base-s}. In forma matematica:

\[
s(t)=
\max
\left[
s_{\min},
\min
\left(
s_{\mathrm{base}},
s_{\mathrm{raw}}(t)
\right)
\right].
\]

Questa limitazione evita valori non realistici: la stubbornness non può superare
il valore massimo iniziale e non può scendere sotto una soglia minima.

\subsection{Pipeline dinamica della simulazione}

\begin{lstlisting}[caption={Procedura principale \texttt{go}}]
;;;;;;;;;;;;;;;;;;;;;;;;;;;;;;;;;;;;;;
;; GO (PIPELINE)
;;;;;;;;;;;;;;;;;;;;;;;;;;;;;;;;;;;;;;

to go

  set model-trading-volume 0

  ;; 1. AGGIORNAMENTO PREZZO E DATI REALI
  if (ticks mod trade-interval = 0) [

    if ticks > 0 [
      aggiorna-dati
    ]

    aggiorna-prezzo-modello
  ]

  ;; 2. AGGIORNAMENTO STATO E MOVIMENTO FISICO
  set current-hype calcola-hype

  ask turtles [
    flock
    repel-different-breeds
    swim
  ]

  ;; 3. TRADING, CORRELAZIONI E PLOT
  if (ticks mod trade-interval = 0) [
    ask turtles [
      compravendita-bitcoin
    ]

    aggiorna-storico-volumi
    aggiorna-correlazione-prezzi

    aggiorna-plot-prezzi-e-volume
    aggiorna-plot-correlazione-volumi
    aggiorna-plot-correlazione-prezzi
  ]

  aggiorna-metriche-stabilita
  aggiorna-plot-parametri

  tick

end
\end{lstlisting}

La procedura \texttt{go} rappresenta il ciclo dinamico principale del modello.
Ogni sua esecuzione corrisponde a un passo temporale della simulazione, ma non
tutte le componenti del modello vengono aggiornate con la stessa frequenza.
Il movimento degli agenti avviene a ogni tick, mentre il mercato, il trading e
alcuni grafici vengono aggiornati solo ogni \texttt{trade-interval} iterazioni.

All'inizio della procedura viene azzerata la variabile

\[
\texttt{model-trading-volume}.
\]

Questa variabile misura il volume di scambio generato dagli agenti nel tick
corrente. Azzerarla all'inizio di ogni ciclo permette di registrare soltanto le
operazioni avvenute nel nuovo passo temporale, evitando che il volume si accumuli
indefinitamente nel tempo.

La prima parte della pipeline riguarda l'aggiornamento dei dati di mercato:

\[
\texttt{if (ticks mod trade-interval = 0) [ ... ]}.
\]

La condizione

\[
\texttt{ticks mod trade-interval = 0}
\]

verifica se il numero corrente di tick è multiplo di \texttt{trade-interval}.
In questo modo, il modello separa la dinamica continua del movimento dalla
dinamica discreta del mercato: gli agenti possono muoversi molte volte prima che
avvenga un nuovo aggiornamento di prezzo e trading.

All'interno di questo blocco, la procedura \texttt{aggiorna-dati} viene chiamata
solo se \texttt{ticks > 0}. Questa scelta evita di aggiornare immediatamente i
dati reali al tick iniziale, dato che il primo prezzo è già stato inizializzato
nella procedura \texttt{setup}. Successivamente viene chiamata

\[
\texttt{aggiorna-prezzo-modello},
\]

che aggiorna il prezzo endogeno generato dal modello. Il prezzo endogeno è il
prezzo prodotto internamente dalla simulazione, a partire dal comportamento degli
agenti, e si distingue dal prezzo reale di Bitcoin, che proviene invece dalla
serie storica esterna.

La seconda parte della procedura aggiorna lo stato comportamentale e il movimento
fisico degli agenti. La riga

\[
\texttt{set current-hype calcola-hype}
\]

calcola il livello corrente di hype o stress di mercato. Questa grandezza dipende
dal rendimento e dalla volatilità del prezzo endogeno. Dal punto di vista
economico, un mercato con alto hype è un mercato più instabile, in cui gli
agenti possono diventare più reattivi e meno conservativi nelle proprie scelte.

Successivamente, tutti gli agenti eseguono tre procedure:

\[
\texttt{flock},
\qquad
\texttt{repel-different-breeds},
\qquad
\texttt{swim}.
\]

La procedura \texttt{flock} aggiorna il comportamento di allineamento collettivo,
facendo sì che gli agenti tendano a coordinarsi con i propri vicini o follower.
La procedura \texttt{repel-different-breeds} introduce invece una repulsione tra
agenti appartenenti a specie diverse, impedendo una sovrapposizione eccessiva tra
gruppi distinti. Infine, \texttt{swim} aggiorna il movimento effettivo degli
agenti nello spazio.

La terza parte della pipeline riguarda trading, correlazioni e grafici. Anche
questa sezione viene eseguita solo ogni \texttt{trade-interval} tick. Gli agenti
eseguono la procedura

\[
\texttt{compravendita-bitcoin},
\]

che determina se comprare o vendere Bitcoin in base alla posizione nello spazio,
alla probabilità di scambio e ai vincoli di portafoglio. Gli agenti non possono
comprare se non hanno liquidità sufficiente e non possono vendere Bitcoin che non
possiedono.

Dopo le operazioni di trading, il modello aggiorna lo storico dei volumi tramite

\[
\texttt{aggiorna-storico-volumi}
\]

e la correlazione tra prezzi reali e prezzi simulati tramite

\[
\texttt{aggiorna-correlazione-prezzi}.
\]

La correlazione di Pearson misura quanto due serie temporali si muovono in modo
linearmente simile. In questo contesto viene usata per confrontare la dinamica
del modello con quella del mercato reale.

Le procedure

\[
\texttt{aggiorna-plot-prezzi-e-volume},
\qquad
\texttt{aggiorna-plot-correlazione-volumi},
\qquad
\texttt{aggiorna-plot-correlazione-prezzi}
\]

aggiornano i grafici relativi all'andamento del prezzo, dei volumi e delle
correlazioni. Poiché questi grafici dipendono da quantità di mercato, vengono
aggiornati solo quando avviene anche l'aggiornamento del trading.

Al di fuori del blocco temporizzato, il modello aggiorna invece le metriche di
stabilità e i parametri collettivi:

\[
\texttt{aggiorna-metriche-stabilita},
\qquad
\texttt{aggiorna-plot-parametri}.
\]

Queste procedure vengono eseguite a ogni tick, perché dipendono dalla dinamica
continua degli agenti e non soltanto dagli eventi di mercato. Possono quindi
monitorare in modo più fine l'evoluzione del flocking, dell'ordine collettivo,
dell'eterogeneità e dell'instabilità.

Infine, il comando

\[
\texttt{tick}
\]

fa avanzare il tempo della simulazione di una unità. La struttura complessiva
della procedura può quindi essere riassunta come

\[
\text{mercato periodico}
\;\longrightarrow\;
\text{hype e movimento}
\;\longrightarrow\;
\text{trading periodico}
\;\longrightarrow\;
\text{metriche e plot}
\;\longrightarrow\;
\text{avanzamento temporale}.
\]

Questa organizzazione consente di rappresentare un sistema in cui gli agenti si
muovono continuamente, mentre le decisioni di mercato e gli aggiornamenti dei
prezzi avvengono con una frequenza più bassa e controllata.

\subsection{Stima del prezzo del modello: fase reale e fase endogena}

\begin{lstlisting}[caption={Procedura di aggiornamento del prezzo del modello}]
;;;;;;;;;;;;;;;;;;;;;;;;;;;;;;;;;;;;;;
;; FUNZIONE DI STIMA DEL PREZZO MODELLO
;;;;;;;;;;;;;;;;;;;;;;;;;;;;;;;;;;;;;;

to aggiorna-prezzo-modello

  ifelse (not current-disabled?) [

    ;; FASE 1: COPIA REALE
    set model-price btc-price
    set model-return btc-return
    set model-volatility btc-volatility

    print (word "[FASE 1 - PREZZO REALE] Giorno: " day-index " | Correlazione Media Volumi: " precision mean-volume-correlation 2)

  ] [

    ;; FASE 2: PREZZO ENDOGENO
    if count turtles > 0 [
      let mean-y mean [ycor] of turtles
      let max-y max-pycor

      let market-pressure (mean-y / max-y)

      let prev-m-price model-price
      set model-price max (list 1 (model-price * (1 + (price-sensitivity * market-pressure))))

      ifelse prev-m-price > 0 [
        set model-return (model-price - prev-m-price) / prev-m-price
        set model-volatility max (list 0.005 (abs model-return))
      ] [
        set model-return 0
        set model-volatility 0.005
      ]
    ]

    print (word ">>> [FASE 2 - PREZZO ENDOGENO] Giorno: " day-index " | Prezzo Stimato: " precision model-price 2)
  ]

  ;; Tracciamento continuo dell'ATH sia in Fase 1 che in Fase 2
  if model-price > btc-ath-price [
    set btc-ath-price model-price
  ]

end
\end{lstlisting}

La procedura \texttt{aggiorna-prezzo-modello} stabilisce come viene determinato
il prezzo utilizzato internamente dalla simulazione. Il codice distingue due
regimi dinamici: una prima fase in cui il modello segue direttamente il mercato
reale e una seconda fase in cui il prezzo viene prodotto endogenamente dagli
agenti.

La condizione

\[
\texttt{ifelse (not current-disabled?) [ ... ] [ ... ]}
\]

seleziona la fase attiva. Se \texttt{current-disabled?} è falso, la corrente
esterna non è disattivata e il modello si trova nella Fase 1. In questa fase il
prezzo del modello viene copiato dal prezzo reale di Bitcoin:

\[
P_m(t)=P_{BTC}(t).
\]

Nel codice questo corrisponde a

\[
\texttt{set model-price btc-price}.
\]

Allo stesso modo, anche rendimento e volatilità del modello vengono posti uguali
a quelli reali:

\[
\texttt{model-return}=\texttt{btc-return},
\qquad
\texttt{model-volatility}=\texttt{btc-volatility}.
\]

Questa fase può essere interpretata come una fase di calibrazione o avvio, in
cui il modello resta agganciato al mercato empirico. Dal punto di vista
finanziario, il rendimento misura la variazione relativa del prezzo tra due
istanti consecutivi:

\[
r(t)=\frac{P(t)-P(t-1)}{P(t-1)}.
\]

La volatilità misura invece l'intensità delle oscillazioni del prezzo. Nel
modello viene usata come grandezza in grado di influenzare la dinamica degli
agenti.

Quando \texttt{current-disabled?} diventa vero, il modello entra nella Fase 2.
In questa fase il prezzo non viene più copiato dal dato reale, ma viene stimato
a partire dallo stato collettivo degli agenti. Il codice calcola innanzitutto la
media delle coordinate verticali:

\[
\texttt{mean-y = mean [ycor] of turtles}.
\]

Nel modello, la coordinata \(y\) ha un significato comportamentale: valori
positivi sono associati a una maggiore propensione all'acquisto, mentre valori
negativi sono associati a una maggiore propensione alla vendita. La media
\(\bar{y}\) fornisce quindi una misura aggregata del sentiment della popolazione.

La pressione di mercato viene definita come

\[
\texttt{market-pressure = mean-y / max-y}.
\]

In forma matematica:

\[
M(t)=\frac{\bar{y}(t)}{y_{\max}}.
\]

Se \(M(t)>0\), gli agenti sono mediamente nella regione di acquisto e il modello
produce una pressione rialzista sul prezzo. Se \(M(t)<0\), gli agenti sono
mediamente nella regione di vendita e il modello produce una pressione ribassista.

Il prezzo endogeno viene aggiornato secondo la regola

\[
P_m(t+1)
=
P_m(t)
\left[
1+\alpha M(t)
\right],
\]

dove \(\alpha\) è rappresentato da \texttt{price-sensitivity}. Nel codice:

\[
\texttt{model-price * (1 + (price-sensitivity * market-pressure))}.
\]

La funzione

\[
\texttt{max (list 1 ...)}
\]

impone un limite inferiore pari a \(1\), evitando che il prezzo stimato diventi
zero o negativo. Questa scelta è coerente con l'interpretazione economica del
prezzo: il prezzo di un'attività finanziaria non dovrebbe assumere valori
negativi.

Dopo l'aggiornamento del prezzo, il codice calcola il rendimento endogeno:

\[
r_m(t)=\frac{P_m(t)-P_m(t-1)}{P_m(t-1)}.
\]

Nel codice:

\[
\texttt{set model-return (model-price - prev-m-price) / prev-m-price}.
\]

La volatilità endogena viene quindi posta uguale al valore assoluto del
rendimento, ma con una soglia minima:

\[
\sigma_m(t)=\max(0.005, |r_m(t)|).
\]

La soglia \(0.005\) rappresenta una volatilità di fondo. Essa evita che la
volatilità diventi esattamente nulla e impedisce che eventuali meccanismi
dipendenti dalla volatilità, come la corrente del modello, si spengano
completamente.

Infine, la procedura aggiorna il massimo prezzo raggiunto:

\[
\texttt{if model-price > btc-ath-price [ set btc-ath-price model-price ]}.
\]

La variabile \texttt{btc-ath-price} conserva quindi il massimo valore osservato
del prezzo usato dal modello, sia durante la fase agganciata al prezzo reale sia
durante la fase endogena. La sigla ATH indica \textit{All-Time High}, cioè il
massimo storico o massimo raggiunto nella finestra di simulazione considerata.

Questa procedura è centrale perché realizza il passaggio da un modello guidato
dai dati reali a un modello autonomo. Nella prima fase il mercato reale fornisce
il riferimento principale; nella seconda fase, invece, il prezzo emerge dallo
stato collettivo degli agenti e dalla loro distribuzione nello spazio
comportamentale.

\subsection{Corrente dinamica nelle due fasi del modello}

\begin{lstlisting}[caption={Componenti della corrente nelle fasi reale ed endogena}]
;;;;;;;;;;;;;;;;;;;;;;;;;;;;;;;;;;;;;;
;; CORRENTE -- COMPONENTE X (DINAMICA FASE 1 / FASE 2)
;;;;;;;;;;;;;;;;;;;;;;;;;;;;;;;;;;;;;;

to-report my-current-x
  let active-vol (ifelse-value (not current-disabled?) [ btc-volatility ] [ model-volatility ])
  report k * active-vol
end


;;;;;;;;;;;;;;;;;;;;;;;;;;;;;;;;;;;;;;
;; CORRENTE -- COMPONENTE Y (DINAMICA FASE 1 / FASE 2)
;;;;;;;;;;;;;;;;;;;;;;;;;;;;;;;;;;;;;;

to-report my-current-y
  let active-ret (ifelse-value (not current-disabled?) [ btc-return ] [ model-return ])

  ifelse breed = agents-a [
    report h * active-ret
  ] [
    ifelse breed = agents-b [
      report (- h * active-ret)
    ] [
      report 0
    ]
  ]
end
\end{lstlisting}

Queste due procedure definiscono la corrente esterna che agisce sugli agenti.
La corrente è divisa in due componenti:

\[
J_x,
\qquad
J_y.
\]

La componente orizzontale è calcolata dal reporter \texttt{my-current-x}, mentre
la componente verticale è calcolata da \texttt{my-current-y}. Entrambe le
componenti dipendono dalla fase dinamica in cui si trova il modello.

La riga

\[
\texttt{let active-vol (ifelse-value (not current-disabled?) [ btc-volatility ] [ model-volatility ])}
\]

seleziona quale volatilità usare. Se \texttt{current-disabled?} è falso, il
modello si trova nella Fase 1 e usa la volatilità reale di Bitcoin,
\texttt{btc-volatility}. Se invece \texttt{current-disabled?} è vero, il modello
si trova nella Fase 2 e usa la volatilità endogena, \texttt{model-volatility}.

La componente orizzontale della corrente è quindi

\[
J_x(t)=k\,\sigma(t),
\]

dove \(k\) è il parametro di intensità della corrente e \(\sigma(t)\) è la
volatilità attiva, reale nella Fase 1 ed endogena nella Fase 2.

Dal punto di vista economico, la volatilità misura l'intensità delle variazioni
di prezzo. Una volatilità elevata descrive un mercato più instabile o turbolento.
Nel modello, questa instabilità viene tradotta in una spinta orizzontale più
forte sugli agenti. Poiché la coordinata \(x\) è collegata alla probabilità di
trading, la corrente orizzontale può essere interpretata come un aumento della
reattività degli agenti durante fasi di mercato più intense.

Il reporter \texttt{my-current-y} definisce invece la componente verticale della
corrente. In questo caso la grandezza rilevante non è la volatilità, ma il
rendimento:

\[
r(t)=\frac{P(t)-P(t-1)}{P(t-1)}.
\]

La riga

\[
\texttt{let active-ret (ifelse-value (not current-disabled?) [ btc-return ] [ model-return ])}
\]

sceglie il rendimento reale nella Fase 1 e il rendimento endogeno nella Fase 2.

La componente verticale dipende anche dalla specie dell'agente. Per gli agenti
della specie \texttt{agents-a}, il codice restituisce

\[
J_y^{(A)}(t)=h\,r(t),
\]

mentre per gli agenti della specie \texttt{agents-b} restituisce

\[
J_y^{(B)}(t)=-h\,r(t).
\]

Per gli agenti della specie \texttt{agents-c}, invece, la componente verticale è
nulla:

\[
J_y^{(C)}(t)=0.
\]

Il parametro \(h\) controlla l'intensità con cui il rendimento influenza il moto
verticale degli agenti. Poiché nel modello la coordinata \(y\) distingue le zone
comportamentali di acquisto e vendita, questa scelta assegna alle specie A e B
una risposta opposta allo stesso segnale di mercato.

Se il rendimento è positivo, gli agenti A vengono spinti verso l'alto, cioè verso
la regione associata all'acquisto, mentre gli agenti B vengono spinti verso il
basso, cioè verso la regione associata alla vendita. Se il rendimento è negativo,
l'effetto si inverte. La specie C, invece, non riceve alcuna spinta verticale
diretta dal rendimento e può essere interpretata come una popolazione più neutra
rispetto alla direzione del mercato.

Questa struttura consente di rappresentare tre categorie comportamentali
differenti: una specie che segue il segnale di rendimento, una specie che reagisce
in modo opposto e una specie neutrale. Il passaggio dalla Fase 1 alla Fase 2
mantiene la stessa forma della corrente, ma cambia la sorgente del segnale:
inizialmente il mercato reale guida la dinamica, successivamente il segnale viene
prodotto dal mercato endogeno generato dagli agenti.

\subsection{Movimento degli agenti con corrente, riflessione ai bordi e rumore}

\begin{lstlisting}[caption={Procedura di movimento degli agenti}]
;;;;;;;;;;;;;;;;;;;;;;;;;;;;;;;;;;;;;;
;; MOVIMENTO
;;;;;;;;;;;;;;;;;;;;;;;;;;;;;;;;;;;;;;

to swim

  set velocity-x velocity-x + my-current-x
  set velocity-y velocity-y + my-current-y

  let next-x xcor + velocity-x
  let next-y ycor + velocity-y

  if next-x > max-pxcor or next-x < min-pxcor [
    set velocity-x (- velocity-x)
    set next-x max (list min-pxcor (min (list max-pxcor next-x)))
  ]

  if next-y > max-pycor or next-y < min-pycor [
    set velocity-y (- velocity-y)
    set next-y max (list min-pycor (min (list max-pycor next-y)))
  ]

  setxy next-x next-y

  let random-angle random-float (2 * temperature) - temperature
  let old-vx velocity-x
  let old-vy velocity-y

  set velocity-x (old-vx * cos random-angle) - (old-vy * sin random-angle)
  set velocity-y (old-vx * sin random-angle) + (old-vy * cos random-angle)

  if velocity-x != 0 or velocity-y != 0 [
    set heading atan velocity-x velocity-y
  ]

end
\end{lstlisting}

La procedura \texttt{swim} aggiorna il movimento degli agenti nello spazio. A
differenza di una dinamica basata soltanto su \texttt{heading} e \texttt{fd}, qui
ogni agente possiede una velocità vettoriale, descritta dalle due componenti

\[
v_x = \texttt{velocity-x},
\qquad
v_y = \texttt{velocity-y}.
\]

La prima parte della procedura aggiunge alla velocità dell'agente la corrente
esterna del modello:

\[
\texttt{set velocity-x velocity-x + my-current-x},
\]

\[
\texttt{set velocity-y velocity-y + my-current-y}.
\]

La corrente rappresenta l'effetto del mercato sul moto degli agenti. La
componente orizzontale dipende dalla volatilità, mentre la componente verticale
dipende dal rendimento e dalla specie dell'agente. In questo modo, l'andamento
del mercato influenza direttamente la traiettoria degli agenti nello spazio
comportamentale.

Dopo l'aggiornamento della velocità, il codice calcola la posizione prevista al
passo successivo:

\[
x_{\mathrm{next}} = x + v_x,
\qquad
y_{\mathrm{next}} = y + v_y.
\]

Nel codice:

\[
\texttt{let next-x xcor + velocity-x},
\qquad
\texttt{let next-y ycor + velocity-y}.
\]

La parte successiva gestisce le pareti riflettenti. Se la posizione prevista
supera il bordo destro o sinistro del mondo, la componente orizzontale della
velocità viene invertita:

\[
v_x \longrightarrow -v_x.
\]

Questo è implementato da

\[
\texttt{set velocity-x (- velocity-x)}.
\]

Successivamente, la posizione viene riportata all'interno dei limiti del mondo
tramite

\[
\texttt{max (list min-pxcor (min (list max-pxcor next-x)))}.
\]

Questa istruzione forza \texttt{next-x} a rimanere compreso tra
\texttt{min-pxcor} e \texttt{max-pxcor}. Lo stesso procedimento viene applicato
alla componente verticale: se l'agente supera il bordo superiore o inferiore, la
velocità verticale viene invertita,

\[
v_y \longrightarrow -v_y,
\]

e la coordinata \texttt{next-y} viene riportata dentro i limiti ammessi.

Il comando

\[
\texttt{setxy next-x next-y}
\]

aggiorna quindi la posizione effettiva dell'agente.

La seconda parte della procedura introduce una perturbazione casuale della
direzione di moto. Viene generato un angolo casuale:

\[
\texttt{random-angle}
\in
[-\texttt{temperature}, \texttt{temperature}).
\]

Il parametro \texttt{temperature} controlla quindi l'intensità del rumore
angolare. Valori bassi producono traiettorie più regolari; valori elevati
rendono il moto più disordinato.

La velocità viene poi ruotata usando la matrice di rotazione bidimensionale:

\[
\begin{pmatrix}
v_x' \\
v_y'
\end{pmatrix}
=
\begin{pmatrix}
\cos\theta & -\sin\theta \\
\sin\theta & \cos\theta
\end{pmatrix}
\begin{pmatrix}
v_x \\
v_y
\end{pmatrix}.
\]

Nel codice, questa rotazione è implementata dalle righe

\[
\texttt{set velocity-x (old-vx * cos random-angle) - (old-vy * sin random-angle)}
\]

e

\[
\texttt{set velocity-y (old-vx * sin random-angle) + (old-vy * cos random-angle)}.
\]

Infine, se la velocità non è nulla, il valore di \texttt{heading} viene
aggiornato in modo coerente con la direzione del vettore velocità:

\[
\texttt{set heading atan velocity-x velocity-y}.
\]

In questo modo, la rappresentazione grafica dell'agente rimane allineata alla
direzione effettiva del suo moto.

Dal punto di vista modellistico, questa procedura combina tre effetti: la corrente
generata dal mercato, la riflessione ai bordi del dominio e il rumore casuale.
Il risultato è una dinamica in cui gli agenti sono spinti dal segnale finanziario,
rimangono confinati nello spazio e mantengono una componente stocastica che può
ostacolare o modificare l'emergere del flocking.

\subsection{Probabilità di compravendita}

\begin{lstlisting}[caption={Reporter per la probabilità di scambio}]
;;;;;;;;;;;;;;;;;;;;;;;;;;;;;;;;;;;;;;
;; PROBABILITA' DI COMPRAVENDITA
;;;;;;;;;;;;;;;;;;;;;;;;;;;;;;;;;;;;;;

to-report probabilita-scambio

  let ampiezza-x (max-pxcor - min-pxcor)

  if ampiezza-x = 0 [ report 0 ]

  let p ((xcor - min-pxcor) / ampiezza-x)

  report min (list 1 (max (list 0 p)))

end
\end{lstlisting}

Il reporter \texttt{probabilita-scambio} assegna a ogni agente una probabilità
di effettuare una compravendita in base alla sua posizione orizzontale nel mondo
NetLogo.

La variabile

\[
\texttt{ampiezza-x}
=
\texttt{max-pxcor} - \texttt{min-pxcor}
\]

rappresenta la larghezza totale del dominio lungo l'asse \(x\). Se tale ampiezza
è nulla, il reporter restituisce direttamente \(0\), evitando una divisione per
zero.

La probabilità grezza viene poi calcolata come

\[
p =
\frac{x - x_{\min}}{x_{\max} - x_{\min}}.
\]

Nel codice:

\[
\texttt{let p ((xcor - min-pxcor) / ampiezza-x)}.
\]

In questo modo, un agente vicino al bordo sinistro del mondo ha una probabilità
di scambio prossima a \(0\), mentre un agente vicino al bordo destro ha una
probabilità prossima a \(1\). La coordinata \(x\) viene quindi interpretata come
una misura della propensione operativa dell'agente: più l'agente si trova verso
destra, maggiore è la probabilità che partecipi al mercato.

L'ultima riga,

\[
\texttt{report min (list 1 (max (list 0 p)))},
\]

serve a vincolare il valore finale nell'intervallo

\[
0 \leq p \leq 1.
\]

Questa operazione evita che eventuali valori numerici anomali producano
probabilità negative o maggiori di uno. Il reporter restituisce quindi sempre
un valore valido interpretabile come probabilità di compravendita.

\subsection{Compravendita di Bitcoin}

\begin{lstlisting}[caption={Procedura di compravendita di Bitcoin degli agenti}]
;;;;;;;;;;;;;;;;;;;;;;;;;;;;;;;;;;;;;;
;; COMPRAVENDITA DI BITCOIN
;;;;;;;;;;;;;;;;;;;;;;;;;;;;;;;;;;;;;;

to compravendita-bitcoin

  let current-p model-price

  if current-p > 0 [

    let p probabilita-scambio

    let quantita-acquisto min (list trade-size (floor (cash-balance / current-p)))
    let quantita-vendita min (list trade-size bitcoin-held)

    if ycor >= 0 [
      if (random-float 1 < p) and (quantita-acquisto > 0) [
        let costo quantita-acquisto * current-p
        set cash-balance cash-balance - costo
        set bitcoin-held bitcoin-held + quantita-acquisto

        set model-trading-volume model-trading-volume + quantita-acquisto
      ]
    ]

    if ycor < 0 [
      if (random-float 1 < p) and (quantita-vendita > 0) [
        let ricavo quantita-vendita * current-p
        set bitcoin-held bitcoin-held - quantita-vendita
        set cash-balance cash-balance + ricavo

        set model-trading-volume model-trading-volume + quantita-vendita
      ]
    ]
  ]

end
\end{lstlisting}

La procedura \texttt{compravendita-bitcoin} descrive il comportamento di trading
degli agenti. Ogni agente può comprare o vendere Bitcoin usando il prezzo
endogeno del modello, memorizzato nella variabile

\[
\texttt{current-p} = \texttt{model-price}.
\]

Il blocco principale viene eseguito solo se il prezzo è positivo:

\[
\texttt{if current-p > 0 [ ... ]}.
\]

Questa condizione evita operazioni economicamente non valide, come acquisti o
vendite a prezzo nullo o negativo.

La probabilità che l'agente effettui uno scambio è data da

\[
\texttt{p = probabilita-scambio}.
\]

Questa probabilità dipende dalla posizione orizzontale dell'agente: più l'agente
si trova verso destra, maggiore è la probabilità che partecipi al mercato.

La quantità massima acquistabile è calcolata come

\[
\texttt{quantita-acquisto}
=
\min
\left(
\texttt{trade-size},
\left\lfloor
\frac{\texttt{cash-balance}}{\texttt{current-p}}
\right\rfloor
\right).
\]

Questa formula impedisce all'agente di comprare più Bitcoin di quanto possa
permettersi con la liquidità disponibile. La funzione \texttt{floor} arrotonda
per difetto il numero di Bitcoin acquistabili.

La quantità massima vendibile è invece

\[
\texttt{quantita-vendita}
=
\min
(
\texttt{trade-size},
\texttt{bitcoin-held}
).
\]

In questo modo l'agente non può vendere più Bitcoin di quelli effettivamente
posseduti.

La coordinata verticale \texttt{ycor} stabilisce il tipo di operazione. Se

\[
\texttt{ycor >= 0},
\]

l'agente si trova nella regione associata all'acquisto. In questo caso, se un
numero casuale tra \(0\) e \(1\) è minore della probabilità \texttt{p} e se la
quantità acquistabile è positiva, l'agente compra Bitcoin. Il costo
dell'operazione è

\[
\texttt{costo}
=
\texttt{quantita-acquisto}
\cdot
\texttt{current-p}.
\]

Il saldo liquido diminuisce di tale costo, mentre la quantità di Bitcoin detenuta
aumenta della quantità acquistata.

Se invece

\[
\texttt{ycor < 0},
\]

l'agente si trova nella regione associata alla vendita. Se la vendita viene
estratta casualmente e l'agente possiede Bitcoin sufficienti, il ricavo è

\[
\texttt{ricavo}
=
\texttt{quantita-vendita}
\cdot
\texttt{current-p}.
\]

In questo caso la quantità di Bitcoin detenuta diminuisce, mentre la liquidità
aumenta del ricavo ottenuto.

Dopo ogni acquisto o vendita riuscita, il volume di trading generato dal modello
viene aggiornato:

\[
\texttt{model-trading-volume}
\leftarrow
\texttt{model-trading-volume}
+
\texttt{quantita}.
\]

Questa variabile misura il numero complessivo di Bitcoin scambiati dagli agenti
nel tick corrente. La procedura collega quindi la posizione spaziale degli agenti
alla loro attività di mercato: la coordinata \(y\) determina la direzione
dell'operazione, mentre la coordinata \(x\) ne influenza la probabilità.

\subsection{Fattore di isteresi dei volumi basato sul prezzo}

\begin{lstlisting}[caption={Reporter per il fattore di scaling dei volumi}]
;;;;;;;;;;;;;;;;;;;;;;;;;;;;;;;;;;;;;;
;; FATTORE DI SCALING VOLUMI (UNIFICATO BASATO SU PREZZO)
;;;;;;;;;;;;;;;;;;;;;;;;;;;;;;;;;;;;;;

to-report volume-hysteresis-factor

  let initial-p max (list 1 btc-initial-price)
  let current-p max (list 1 model-price)
  let ath-p     max (list current-p btc-ath-price)

  let ratio-current (current-p / initial-p)
  let ratio-ath     (ath-p / initial-p)

  ;; Formula continua dell'isteresi valida sia in Fase 1 che in Fase 2
  let price-hysteresis-state ratio-current + (beta-vol * (ratio-ath - ratio-current))

  report max (list 0.01 price-hysteresis-state)

end
\end{lstlisting}

Il reporter \texttt{volume-hysteresis-factor} calcola un fattore di scaling dei
volumi a partire dall'andamento del prezzo. L'idea è che il volume di mercato non
dipenda soltanto dal prezzo corrente, ma anche dalla memoria del massimo prezzo
raggiunto in passato.

Le prime tre variabili definiscono prezzi sempre positivi:

\[
\texttt{initial-p} = \max(1,\texttt{btc-initial-price}),
\]

\[
\texttt{current-p} = \max(1,\texttt{model-price}),
\]

\[
\texttt{ath-p} = \max(\texttt{current-p},\texttt{btc-ath-price}).
\]

L'uso di \(\max(1,\cdot)\) evita divisioni per zero o valori troppo piccoli. La
variabile \texttt{ath-p} rappresenta invece il massimo prezzo rilevante per il
modello, cioè il maggiore tra il prezzo corrente e il massimo storico registrato.

Il codice calcola poi due rapporti normalizzati rispetto al prezzo iniziale:

\[
\texttt{ratio-current}
=
\frac{\texttt{current-p}}{\texttt{initial-p}},
\]

\[
\texttt{ratio-ath}
=
\frac{\texttt{ath-p}}{\texttt{initial-p}}.
\]

Il primo rapporto misura quanto il prezzo corrente è cresciuto rispetto al
valore iniziale. Il secondo misura invece quanto il massimo storico è cresciuto
rispetto al valore iniziale.

Il fattore di isteresi viene definito come

\[
H_V
=
\texttt{ratio-current}
+
\beta_{\mathrm{vol}}
\left(
\texttt{ratio-ath}
-
\texttt{ratio-current}
\right).
\]

Nel codice:

\[
\texttt{price-hysteresis-state = ratio-current + (beta-vol * (ratio-ath - ratio-current))}.
\]

Il parametro \texttt{beta-vol} controlla quanto peso viene dato alla memoria del
massimo storico. Se \texttt{beta-vol} è vicino a \(0\), il fattore dipende quasi
solo dal prezzo corrente. Se invece è più alto, il fattore conserva maggiormente
l'effetto dei picchi di prezzo passati.

In questo senso si parla di isteresi: lo stato attuale del volume non dipende
solo dalla situazione presente, ma anche dalla storia precedente del prezzo.
Anche se il prezzo scende dopo aver raggiunto un massimo, il modello può
mantenere un livello di volume più elevato a causa della memoria del picco.

L'ultima riga

\[
\texttt{report max (list 0.01 price-hysteresis-state)}
\]

impone un limite inferiore pari a \(0.01\). In questo modo il fattore di scaling
dei volumi resta sempre positivo e non può annullarsi completamente.

\subsection{Volume del modello espresso in dollari}

\begin{lstlisting}[caption={Reporter per il volume di trading del modello in dollari}]
;;;;;;;;;;;;;;;;;;;;;;;;;;;;;;;;;;;;;;
;; VOLUME DEL MODELLO IN DOLLARI
;;;;;;;;;;;;;;;;;;;;;;;;;;;;;;;;;;;;;;

to-report model-trading-volume-usd

  let current-p model-price
  if current-p <= 0 [ report 0 ]

  report (model-trading-volume * current-p) * volume-hysteresis-factor

end
\end{lstlisting}

Il reporter \texttt{model-trading-volume-usd} converte il volume di trading
generato dagli agenti in un controvalore espresso in dollari.

La variabile locale

\[
\texttt{current-p} = \texttt{model-price}
\]

memorizza il prezzo endogeno corrente del modello. Questo prezzo viene usato per
trasformare il numero di Bitcoin scambiati in un valore monetario.

Il controllo

\[
\texttt{if current-p <= 0 [ report 0 ]}
\]

serve a evitare che il volume venga calcolato usando un prezzo nullo o negativo.
In tal caso il reporter restituisce direttamente \(0\), perché non avrebbe senso
economico calcolare un controvalore di mercato con un prezzo non positivo.

Se il prezzo è valido, il volume in dollari viene calcolato come

\[
V_{\mathrm{model}}^{USD}
=
V_{\mathrm{model}}^{BTC}
\cdot
P_m
\cdot
H_V,
\]

dove \(V_{\mathrm{model}}^{BTC}\) è il volume di Bitcoin scambiato dagli agenti,
\(P_m\) è il prezzo endogeno del modello e \(H_V\) è il fattore di isteresi dei
volumi.

Nel codice questa formula corrisponde a

\[
\texttt{(model-trading-volume * current-p) * volume-hysteresis-factor}.
\]

Il termine \texttt{model-trading-volume * current-p} trasforma il volume da
Bitcoin a dollari. Il fattore \texttt{volume-hysteresis-factor} introduce invece
una correzione legata alla memoria del prezzo: i volumi non dipendono soltanto
dagli scambi correnti, ma anche dalla dinamica storica del prezzo e dal massimo
raggiunto in precedenza.

In questo modo il modello produce un volume simulato in dollari confrontabile con
il volume reale di mercato.

\subsection{Lettura dei dati storici da file CSV}

\begin{lstlisting}[caption={Procedura per il caricamento dei dati da un file CSV}]
;;;;;;;;;;;;;;;;;;;;;;;;;;;;;;;;;;;;;;
;; LETTURA FILE CSV
;;;;;;;;;;;;;;;;;;;;;;;;;;;;;;;;;;;;;;

to carica-dati-da-csv [ nome-file ]
  set prices []
  file-close-all

  carefully [
    let dati csv:from-file nome-file

    if length dati > 0 [
      let prima-riga item 0 dati
      if not is-number? item 1 prima-riga [
        set dati but-first dati
      ]
    ]

    let t 0
    foreach dati [ riga ->
      if length riga >= 2 [
        let p item 1 riga
        if is-number? p [
          let vol 0
          if length riga >= 3 [
            if is-number? item 2 riga [
              set vol item 2 riga
            ]
          ]
          set prices lput (list t p vol) prices
          set t t + 1
        ]
      ]
    ]

    print (word "CSV caricato con successo! Estratti " length prices " record storici validi.")
  ] [
    file-close-all
    user-message (word "Errore nella lettura del file CSV: " nome-file "\nVerifica che il file sia nella stessa cartella.")
  ]
end
\end{lstlisting}

La procedura \texttt{carica-dati-da-csv} permette di caricare una serie storica
da un file CSV esterno. L'argomento \texttt{nome-file} indica il nome del file da
leggere, che deve trovarsi nella cartella del modello NetLogo.

All'inizio la lista \texttt{prices} viene svuotata:

\[
\texttt{set prices []}.
\]

In questo modo, ogni nuovo caricamento sostituisce i dati precedenti invece di
aggiungersi a essi. Il comando \texttt{file-close-all} chiude eventuali file
rimasti aperti da letture precedenti.

La lettura vera e propria avviene dentro un blocco \texttt{carefully}. Questo
serve a intercettare eventuali errori, ad esempio se il file non esiste, se il
nome è scritto male oppure se il file non è nella cartella corretta.

Il comando

\[
\texttt{let dati csv:from-file nome-file}
\]

legge il file CSV e lo trasforma in una lista di righe. Ogni riga è a sua volta
una lista contenente i valori delle colonne.

Il blocco successivo controlla se la prima riga è un'intestazione. Se nella
seconda colonna della prima riga non è presente un numero, allora il codice
interpreta quella riga come header e la rimuove:

\[
\texttt{set dati but-first dati}.
\]

Successivamente viene introdotta la variabile \texttt{t}, che rappresenta
l'indice temporale dei dati caricati. Per ogni riga del file, il codice controlla
che siano presenti almeno due colonne. La seconda colonna viene interpretata come
prezzo:

\[
\texttt{let p item 1 riga}.
\]

Se \texttt{p} è un numero, allora la riga viene considerata valida. Il volume
viene inizializzato a zero:

\[
\texttt{let vol 0}.
\]

Se esiste anche una terza colonna e questa contiene un valore numerico, allora
quel valore viene usato come volume:

\[
\texttt{set vol item 2 riga}.
\]

Ogni record valido viene infine inserito nella lista \texttt{prices} nella forma

\[
(t, p, vol),
\]

cioè:

\[
\texttt{list t p vol}.
\]

Nel modello, quindi, ogni elemento di \texttt{prices} contiene tre informazioni:
l'indice temporale, il prezzo e il volume associato. Dopo ogni riga valida,
\texttt{t} viene incrementato di uno, così da costruire una sequenza temporale
ordinata.

Se la lettura va a buon fine, il modello stampa un messaggio con il numero di
record storici validi estratti. Se invece si verifica un errore, viene chiuso ogni
file aperto e compare un messaggio per avvisare l'utente che il CSV non è stato
letto correttamente.

\subsection{Gestione degli scenari di simulazione}

\begin{lstlisting}[caption={Procedura per la selezione dello scenario di mercato}]
;;;;;;;;;;;;;;;;;;;;;;;;;;;;;;;;;;;;;;
;; GESTIONE SCENARI
;;;;;;;;;;;;;;;;;;;;;;;;;;;;;;;;;;;;;;

to scarica-dati

  if scenario = "csv" [
    carica-dati-da-csv "dati_btc.csv"
    stop
  ]

  if scenario = "api-live" [
    set prices genera-prezzi-api
    stop
  ]

  if scenario = "crash" [
    set prices genera-crash 60000
    stop
  ]

  if scenario = "bull-market" [
    set prices genera-prezzi 30000 0.015 0.01
    stop
  ]

  if scenario = "bear-market" [
    set prices genera-prezzi 60000 -0.015 0.01
    stop
  ]

  if scenario = "sideways" [
    set prices genera-prezzi 40000 0 0.005
    stop
  ]

  if scenario = "bubble" [
    set prices genera-bolla 30000
    stop
  ]

end
\end{lstlisting}

La procedura \texttt{scarica-dati} seleziona quale serie storica utilizzare per
inizializzare il modello. La scelta dipende dal valore della variabile
\texttt{scenario}, che può corrispondere a diversi tipi di mercato.

Se lo scenario è

\[
\texttt{"csv"},
\]

il modello carica i dati dal file esterno \texttt{dati\_btc.csv} tramite la
procedura

\[
\texttt{carica-dati-da-csv "dati\_btc.csv"}.
\]

In questo caso la simulazione usa una serie storica salvata localmente.

Se lo scenario è

\[
\texttt{"api-live"},
\]

la lista \texttt{prices} viene costruita tramite

\[
\texttt{set prices genera-prezzi-api}.
\]

Questo scenario permette di usare dati ottenuti da una sorgente esterna, invece
di dati generati artificialmente.

Gli altri scenari producono invece serie sintetiche. Lo scenario

\[
\texttt{"crash"}
\]

genera un andamento caratterizzato da una forte caduta del prezzo, partendo dal
valore iniziale \(60000\). Lo scenario

\[
\texttt{"bull-market"}
\]

genera una dinamica crescente, con prezzo iniziale \(30000\), drift positivo
\(0.015\) e rumore \(0.01\). Lo scenario

\[
\texttt{"bear-market"}
\]

usa invece un drift negativo pari a \(-0.015\), simulando un mercato
tendenzialmente ribassista.

Lo scenario

\[
\texttt{"sideways"}
\]

rappresenta un mercato laterale, cioè privo di una direzione dominante. Infatti
il drift è nullo:

\[
\texttt{genera-prezzi 40000 0 0.005}.
\]

Infine, lo scenario

\[
\texttt{"bubble"}
\]

genera una dinamica di bolla, partendo da un prezzo iniziale pari a \(30000\).

Dopo ogni scelta viene eseguito il comando

\[
\texttt{stop}.
\]

Questo interrompe la procedura appena uno scenario è stato riconosciuto, evitando
che vengano controllati inutilmente anche gli scenari successivi.

La procedura permette quindi di usare lo stesso modello in condizioni di mercato
diverse: dati reali da CSV, dati esterni tramite API oppure scenari sintetici
come crescita, crisi, mercato laterale e bolla speculativa.

\subsection{Generazione dei prezzi tramite API live}

\begin{lstlisting}[caption={Reporter per la costruzione di una serie prezzi con prezzo finale live da API}]
;;;;;;;;;;;;;;;;;;;;;;;;;;;;;;;;;;;;;;
;; PREZZO LIVE DA API
;;;;;;;;;;;;;;;;;;;;;;;;;;;;;;;;;;;;;;

to-report genera-prezzi-api

  let url "https://api.coinbase.com/v2/prices/BTC-USD/spot"
  let raw-json ""
  let fetch-failed? false

  carefully [
    set raw-json fetch:url url
  ] [
    set fetch-failed? true
  ]

  if fetch-failed? or raw-json = "" [
    user-message "Errore richiesta API\nUso del trend sintetico di fallback."
    report genera-prezzi 40000 0 0.005
  ]

  let data table:from-json raw-json
  let data-map table:get data "data"
  let current-price read-from-string (table:get data-map "amount")

  let result []
  let p 40000
  let t 0

  while [t < (num-days - 1)] [
    let r (random-normal 0 0.005)
    set p p * (1 + r)
    set result lput (list t p 0) result
    set t t + 1
  ]

  set result lput (list t current-price 0) result

  report result

end
\end{lstlisting}

Il reporter \texttt{genera-prezzi-api} costruisce una serie di prezzi usando un
prezzo live ottenuto tramite API. L'obiettivo non è scaricare un'intera serie
storica reale, ma generare una traiettoria sintetica che termina con il prezzo
corrente di Bitcoin.

La variabile

\[
\texttt{url}
\]

contiene l'indirizzo dell'API da cui viene richiesto il prezzo spot di
\texttt{BTC-USD}. Le variabili \texttt{raw-json} e \texttt{fetch-failed?}
servono rispettivamente a memorizzare la risposta grezza dell'API e a registrare
l'eventuale fallimento della richiesta.

La chiamata esterna viene eseguita dentro un blocco \texttt{carefully}:

\[
\texttt{set raw-json fetch:url url}.
\]

Se la richiesta fallisce, la variabile \texttt{fetch-failed?} viene posta uguale
a \texttt{true}. In questo modo il modello non si blocca in caso di assenza di
connessione, errore dell'API o risposta non valida.

Se la richiesta non va a buon fine, oppure se la risposta è vuota, il modello
mostra un messaggio all'utente e usa una serie sintetica di fallback:

\[
\texttt{genera-prezzi 40000 0 0.005}.
\]

Questo fallback rappresenta un mercato laterale con prezzo iniziale \(40000\),
drift nullo e rumore pari a \(0.005\).

Se invece la risposta dell'API è valida, il testo JSON viene convertito in una
struttura leggibile tramite

\[
\texttt{table:from-json raw-json}.
\]

Il codice accede poi al campo \texttt{"data"} e recupera il valore
\texttt{"amount"}, che contiene il prezzo corrente. Poiché tale valore arriva
come stringa, viene trasformato in numero con

\[
\texttt{read-from-string}.
\]

La seconda parte del reporter costruisce la serie \texttt{result}. La variabile
\texttt{p} parte da \(40000\), mentre \texttt{t} rappresenta l'indice temporale.
Finché

\[
t < \texttt{num-days} - 1,
\]

il modello genera un rendimento casuale

\[
r \sim \mathcal{N}(0,0.005)
\]

e aggiorna il prezzo secondo la regola

\[
p_{t+1}=p_t(1+r).
\]

Ogni valore generato viene aggiunto alla lista \texttt{result} nella forma

\[
(t,p,0),
\]

dove il terzo elemento rappresenta il volume, qui posto uguale a zero.

Alla fine del ciclo, l'ultimo elemento della serie viene sostituito dal prezzo
live ottenuto tramite API:

\[
\texttt{set result lput (list t current-price 0) result}.
\]

Il reporter restituisce quindi una lista di record temporali in cui i primi
valori sono generati sinteticamente, mentre l'ultimo prezzo coincide con il
prezzo corrente letto dall'API.

\subsection{Generazione di scenari sintetici di prezzo}

\begin{lstlisting}[caption={Reporter per la generazione di prezzi sintetici, crash e bolla}]
;;;;;;;;;;;;;;;;;;;;;;;;;;;;;;;;;;;;;;
;; GENERA PREZZI SINTETICI
;;;;;;;;;;;;;;;;;;;;;;;;;;;;;;;;;;;;;;

to-report genera-prezzi [ prezzo-iniziale drift sigma ]

  let result []
  let p prezzo-iniziale
  let t 0

  while [t < num-days] [
    let r drift + random-normal 0 sigma
    set p p * (1 + r)
    if p < 1 [ set p 1 ]
    set result lput (list t p 0) result
    set t t + 1
  ]

  report result

end


;;;;;;;;;;;;;;;;;;;;;;;;;;;;;;;;;;;;;;
;; GENERA CRASH SINTETICO
;;;;;;;;;;;;;;;;;;;;;;;;;;;;;;;;;;;;;;

to-report genera-crash [prezzo-iniziale]

  let result []
  let p prezzo-iniziale
  let t 0

  while [t < num-days] [
    let r 0

    ifelse t < 100 [
      set r random-normal 0 0.005
    ][
      ifelse t < 110 [
        set r -0.25 + random-normal 0 0.02
      ][
        set r random-normal 0.003 0.012
      ]
    ]

    set p p * (1 + r)
    if p < 1 [ set p 1 ]
    set result lput (list t p 0) result
    set t t + 1
  ]

  report result

end


;;;;;;;;;;;;;;;;;;;;;;;;;;;;;;;;;;;;;;
;; GENERA BOLLA SINTETICA
;;;;;;;;;;;;;;;;;;;;;;;;;;;;;;;;;;;;;;

to-report genera-bolla [prezzo-iniziale]

  let result []
  let p prezzo-iniziale
  let t 0

  while [t < num-days] [
    let r 0

    ifelse t < 100 [
      set r 0.025 + random-normal 0 0.008
    ][
      ifelse t < 110 [
        set r 0.08 + random-normal 0 0.02
      ][
        ifelse t < 115 [
          set r -0.3 + random-normal 0 0.025
        ][
          set r -0.002 + random-normal 0 0.015
        ]
      ]
    ]

    set p p * (1 + r)
    if p < 1 [ set p 1 ]
    set result lput (list t p 0) result
    set t t + 1
  ]

  report result

end
\end{lstlisting}

Questi tre reporter generano serie temporali sintetiche di prezzo, usate quando
non si vogliono caricare dati reali o quando si desidera studiare il comportamento
del modello in scenari controllati.

Il reporter \texttt{genera-prezzi} produce una dinamica generica a partire da tre
parametri: \texttt{prezzo-iniziale}, \texttt{drift} e \texttt{sigma}. Il prezzo
parte dal valore iniziale assegnato e, a ogni passo temporale, viene aggiornato
tramite un rendimento casuale:

\[
r = \texttt{drift} + \mathcal{N}(0,\texttt{sigma}).
\]

Il parametro \texttt{drift} rappresenta la tendenza media del mercato: se è
positivo genera una crescita media, se è negativo genera una discesa media, se è
nullo produce un mercato laterale. Il parametro \texttt{sigma} controlla invece
l'intensità delle oscillazioni casuali.

Il prezzo viene aggiornato secondo la regola

\[
p_{t+1}=p_t(1+r).
\]

La condizione

\[
\texttt{if p < 1 [ set p 1 ]}
\]

impone un limite inferiore al prezzo, evitando che la serie assuma valori nulli o
negativi. Ogni osservazione viene salvata nella lista \texttt{result} nella forma

\[
(t,p,0),
\]

dove \(t\) è l'indice temporale, \(p\) è il prezzo e il terzo valore rappresenta
il volume, qui posto uguale a zero.

Il reporter \texttt{genera-crash} costruisce invece uno scenario di crollo del
mercato. La dinamica è divisa in tre fasi. Per i primi \(100\) tick il prezzo
oscilla intorno a un rendimento medio nullo:

\[
r \sim \mathcal{N}(0,0.005).
\]

Tra il tick \(100\) e il tick \(109\), il modello introduce una forte caduta:

\[
r = -0.25 + \mathcal{N}(0,0.02).
\]

Dopo il crash, il rendimento torna moderatamente positivo:

\[
r \sim \mathcal{N}(0.003,0.012).
\]

Questo permette di simulare una fase inizialmente stabile, seguita da un crollo
rapido e da una successiva fase di parziale recupero.

Il reporter \texttt{genera-bolla} genera invece uno scenario di bolla
speculativa. Anche in questo caso la dinamica è divisa in più fasi. Nei primi
\(100\) tick il prezzo cresce in modo sostenuto:

\[
r = 0.025 + \mathcal{N}(0,0.008).
\]

Tra il tick \(100\) e il tick \(109\), la crescita accelera ulteriormente:

\[
r = 0.08 + \mathcal{N}(0,0.02).
\]

Tra il tick \(110\) e il tick \(114\), la bolla esplode e il prezzo subisce una
brusca caduta:

\[
r = -0.3 + \mathcal{N}(0,0.025).
\]

Infine, dopo il collasso, il mercato entra in una fase discendente e più
rumorosa:

\[
r = -0.002 + \mathcal{N}(0,0.015).
\]

Queste tre procedure permettono quindi di confrontare il comportamento degli
agenti in mercati con caratteristiche molto diverse: crescita regolare, ribasso,
stabilità laterale, crash improvviso e bolla speculativa.

\subsection{Aggiornamento giornaliero dei dati di mercato}

\begin{lstlisting}[caption={Procedura di aggiornamento giornaliero dei prezzi e dei volumi reali}]
;;;;;;;;;;;;;;;;;;;;;;;;;;;;;;;;;;;;;;
;; AGGIORNAMENTO GIORNALIERO
;;;;;;;;;;;;;;;;;;;;;;;;;;;;;;;;;;;;;;

to aggiorna-dati

  if not is-list? prices [ stop ]
  if day-index >= length prices [ stop ]

  let record item day-index prices
  let new-price item 1 record

  ifelse length record >= 3 [
    set real-trading-volume-usd item 2 record
  ] [
    set real-trading-volume-usd 0
  ]

  if real-volume-initial = 0 and real-trading-volume-usd > 0 [
    set real-volume-initial real-trading-volume-usd
    set real-volume-ath real-trading-volume-usd
  ]

  if real-trading-volume-usd > real-volume-ath [
    set real-volume-ath real-trading-volume-usd
  ]

  ifelse btc-prev-price > 0 [
    set btc-return (new-price - btc-prev-price) / btc-prev-price
    set btc-volatility abs btc-return
  ][
    set btc-return 0
    set btc-volatility 0
  ]

  set btc-prev-price btc-price
  set btc-price new-price

  if btc-initial-price = 0 and btc-price > 0 [
    set btc-initial-price btc-price
    set btc-ath-price btc-price
  ]

  if btc-price > btc-ath-price [
    set btc-ath-price btc-price
  ]

  set day-index day-index + 1

end
\end{lstlisting}

La procedura \texttt{aggiorna-dati} aggiorna i dati reali di mercato letti dalla
lista \texttt{prices}. Ogni elemento della lista rappresenta un record temporale
contenente almeno il prezzo e, se disponibile, anche il volume reale di scambio.

All'inizio vengono effettuati due controlli di sicurezza:

\[
\texttt{if not is-list? prices [ stop ]}
\]

e

\[
\texttt{if day-index >= length prices [ stop ]}.
\]

Il primo controllo verifica che \texttt{prices} sia effettivamente una lista; il
secondo evita di leggere oltre la fine della serie storica disponibile.

Il record corrente viene estratto tramite

\[
\texttt{let record item day-index prices},
\]

mentre il nuovo prezzo viene preso dalla seconda posizione del record:

\[
\texttt{let new-price item 1 record}.
\]

Se il record contiene almeno tre elementi, il terzo valore viene interpretato
come volume reale in dollari:

\[
\texttt{set real-trading-volume-usd item 2 record}.
\]

Se invece il volume non è disponibile, viene posto uguale a zero. Questa scelta
permette alla procedura di funzionare sia con serie storiche complete, contenenti
prezzo e volume, sia con serie contenenti soltanto il prezzo.

Il blocco successivo inizializza il volume reale iniziale e il massimo volume
raggiunto:

\[
\texttt{real-volume-initial},
\qquad
\texttt{real-volume-ath}.
\]

Il volume iniziale viene impostato solo quando viene osservato il primo volume
positivo. Successivamente, se il volume reale corrente supera il massimo
registrato, la variabile \texttt{real-volume-ath} viene aggiornata.

La procedura calcola poi il rendimento reale di Bitcoin:

\[
r_{BTC}(t)
=
\frac{P_{BTC}(t)-P_{BTC}(t-1)}
     {P_{BTC}(t-1)}.
\]

Nel codice, questa quantità è memorizzata in \texttt{btc-return}:

\[
\texttt{set btc-return (new-price - btc-prev-price) / btc-prev-price}.
\]

La volatilità reale viene invece approssimata tramite il valore assoluto del
rendimento:

\[
\sigma_{BTC}(t)=|r_{BTC}(t)|.
\]

Se il prezzo precedente non è positivo, rendimento e volatilità vengono posti
uguali a zero, evitando divisioni non valide.

Dopo il calcolo del rendimento, il codice aggiorna il prezzo precedente e il
prezzo corrente:

\[
\texttt{set btc-prev-price btc-price},
\qquad
\texttt{set btc-price new-price}.
\]

In questo modo il modello mantiene memoria del prezzo usato nel passo temporale
precedente e registra il nuovo prezzo reale appena letto dalla serie storica.

Se \texttt{btc-initial-price} non è ancora stato assegnato, il primo prezzo
positivo viene usato sia come prezzo iniziale sia come massimo iniziale:

\[
\texttt{btc-initial-price} = \texttt{btc-price},
\qquad
\texttt{btc-ath-price} = \texttt{btc-price}.
\]

Infine, se il nuovo prezzo supera il massimo storico registrato, viene aggiornato
\texttt{btc-ath-price}. La sigla ATH indica \emph{All-Time High}, cioè il massimo
prezzo raggiunto nella finestra di simulazione.

La procedura termina incrementando \texttt{day-index} di una unità. In questo
modo, alla chiamata successiva, il modello leggerà il record temporale seguente
della lista \texttt{prices}.

\subsection{Correlazione dei volumi e switch automatico di regime}

\begin{lstlisting}[caption={Calcolo della correlazione di Pearson sui volumi e aggiornamento del relativo plot}]
;;;;;;;;;;;;;;;;;;;;;;;;;;;;;;;;;;;;;;
;; CALCOLO E PLOT CORRELAZIONE VOLUMI (PEARSON)
;;;;;;;;;;;;;;;;;;;;;;;;;;;;;;;;;;;;;;

to aggiorna-storico-volumi

  set model-volume-history lput model-trading-volume-usd model-volume-history
  set real-volume-history lput real-trading-volume-usd real-volume-history

  ;; Applicazione finestra mobile agli storici dei volumi
  if is-number? window-size and window-size > 0 [
    if length model-volume-history > window-size [
      set model-volume-history but-first model-volume-history
      set real-volume-history but-first real-volume-history
    ]
  ]

  ;; Calcolo correlazione di Pearson per i volumi
  if length model-volume-history >= 2 [
    let mean-m mean model-volume-history
    let mean-r mean real-volume-history

    let cov sum (map [ [m r] -> (m - mean-m) * (r - mean-r) ] model-volume-history real-volume-history)
    let std-m sqrt sum (map [ m -> (m - mean-m) ^ 2 ] model-volume-history)
    let std-r sqrt sum (map [ r -> (r - mean-r) ^ 2 ] real-volume-history)

    ifelse (std-m > 0 and std-r > 0) [
      set volume-correlation cov / (std-m * std-r)
    ] [
      set volume-correlation 0
    ]

    ;; Applicazione finestra mobile allo storico della correlazione
    set volume-correlation-history lput volume-correlation volume-correlation-history
    if length volume-correlation-history > window-size [
      set volume-correlation-history but-first volume-correlation-history
    ]

    set mean-volume-correlation mean volume-correlation-history
  ]

  if (not current-disabled?) and (length model-volume-history >= window-size) [

    ifelse mean-volume-correlation >= pearson-threshold [
      set above-threshold-counter above-threshold-counter + 1
    ] [
      set above-threshold-counter 0
    ]

    if above-threshold-counter > 5 [
      set current-disabled? true
      print (word "!!! REGIME RAGGIUNTO A GIORNO " day-index " !!!")
      print (word "Correlazione media volumi (" precision mean-volume-correlation 3 ") >= " pearson-threshold " per piu' di 5 tick consecutivi.")
      print "Passaggio a FASE 2 (Prezzo Endogeno)."
    ]
  ]

end

to aggiorna-plot-correlazione-volumi

  carefully [
    set-current-plot "Correlazione Volumi (Pearson r)"

    assicura-penna "soglia-switch" red
    plotxy day-index pearson-threshold

    assicura-penna "zero" black
    plotxy day-index 0

    if real-trading-volume-usd > 0 [
      assicura-penna "correlazione" blue
      plotxy day-index volume-correlation

      assicura-penna "correlazione-media" green
      plotxy day-index mean-volume-correlation
    ]
  ] [ ]

end
\end{lstlisting}

La procedura \texttt{aggiorna-storico-volumi} aggiorna gli storici dei volumi
generati dal modello e dei volumi reali. A ogni chiamata, il valore corrente del
volume simulato in dollari viene aggiunto a \texttt{model-volume-history}, mentre
il volume reale viene aggiunto a \texttt{real-volume-history}.

Per evitare che gli storici crescano indefinitamente, il codice applica una
finestra mobile di ampiezza \texttt{window-size}. Se la lunghezza degli storici
supera tale valore, viene eliminato il primo elemento tramite \texttt{but-first}.
In questo modo, la correlazione viene calcolata solo sugli ultimi valori
osservati, non sull'intera storia della simulazione.

Quando sono disponibili almeno due osservazioni, il codice calcola la
correlazione di Pearson tra la serie dei volumi simulati e la serie dei volumi
reali. Vengono prima calcolate le medie:

\[
\bar{V}_m = \texttt{mean model-volume-history},
\qquad
\bar{V}_r = \texttt{mean real-volume-history}.
\]

Successivamente vengono calcolate la covarianza non normalizzata e le deviazioni
standard non normalizzate:

\[
\text{cov}
=
\sum_i
\left(
V_{m,i}-\bar{V}_m
\right)
\left(
V_{r,i}-\bar{V}_r
\right),
\]

\[
\sigma_m
=
\sqrt{
\sum_i
\left(
V_{m,i}-\bar{V}_m
\right)^2
},
\qquad
\sigma_r
=
\sqrt{
\sum_i
\left(
V_{r,i}-\bar{V}_r
\right)^2
}.
\]

La correlazione viene quindi calcolata come

\[
r_V =
\frac{\text{cov}}{\sigma_m \sigma_r}.
\]

Se almeno una delle due deviazioni standard è nulla, la correlazione viene posta
uguale a \(0\), perché non è possibile misurare una relazione lineare quando una
delle due serie è costante.

Il valore ottenuto viene salvato in \texttt{volume-correlation} e aggiunto allo
storico \texttt{volume-correlation-history}. Anche questo storico viene mantenuto
entro la dimensione massima \texttt{window-size}. La variabile
\texttt{mean-volume-correlation} rappresenta la media recente delle correlazioni
dei volumi.

La parte finale della procedura gestisce lo switch automatico tra Fase 1 e Fase
2. Finché \texttt{current-disabled?} è falso, il modello resta nella fase guidata
dai dati reali. Quando lo storico dei volumi ha raggiunto la lunghezza della
finestra mobile, il codice controlla se

\[
\texttt{mean-volume-correlation} \geq \texttt{pearson-threshold}.
\]

Se la condizione è verificata, il contatore
\texttt{above-threshold-counter} aumenta di \(1\). Se invece la correlazione
media scende sotto soglia, il contatore viene azzerato. Quando il contatore
supera \(5\), il modello considera raggiunto un regime sufficientemente stabile
e imposta

\[
\texttt{current-disabled? = true}.
\]

Questo provoca il passaggio alla Fase 2, in cui il prezzo non viene più copiato
dal mercato reale ma viene generato endogenamente dal modello.

La procedura \texttt{aggiorna-plot-correlazione-volumi} aggiorna il grafico
\texttt{Correlazione Volumi (Pearson r)}. Il blocco \texttt{carefully} evita che
eventuali problemi con il plot o con le penne blocchino la simulazione.

Nel grafico vengono tracciate quattro curve o riferimenti:

\[
\texttt{soglia-switch},
\qquad
\texttt{zero},
\qquad
\texttt{correlazione},
\qquad
\texttt{correlazione-media}.
\]

La penna \texttt{soglia-switch} disegna la soglia
\texttt{pearson-threshold}, mentre la penna \texttt{zero} disegna la linea di
riferimento \(r=0\). Se il volume reale è positivo, vengono poi tracciate la
correlazione istantanea dei volumi e la sua media mobile. In questo modo il plot
mostra sia il valore corrente della correlazione sia la sua tendenza recente,
rendendo visibile il momento in cui il modello si avvicina alla soglia di switch.

\subsection{Correlazione di Pearson tra prezzi reali e prezzi simulati}

\begin{lstlisting}[caption={Calcolo della correlazione di Pearson sui prezzi e aggiornamento del relativo plot}]
;;;;;;;;;;;;;;;;;;;;;;;;;;;;;;;;;;;;;;
;; CALCOLO E PLOT CORRELAZIONE PREZZI (PEARSON)
;;;;;;;;;;;;;;;;;;;;;;;;;;;;;;;;;;;;;;

to aggiorna-correlazione-prezzi

  set model-price-history lput model-price model-price-history
  set real-price-history lput btc-price real-price-history

  ;; Finestra mobile sugli storici dei prezzi
  if is-number? window-size and window-size > 0 [
    if length model-price-history > window-size [
      set model-price-history but-first model-price-history
      set real-price-history but-first real-price-history
    ]
  ]

  ;; Calcolo correlazione di Pearson per i prezzi
  if length model-price-history >= 2 [
    let mean-m mean model-price-history
    let mean-r mean real-price-history

    let cov sum (map [ [m r] -> (m - mean-m) * (r - mean-r) ] model-price-history real-price-history)
    let std-m sqrt sum (map [ m -> (m - mean-m) ^ 2 ] model-price-history)
    let std-r sqrt sum (map [ r -> (r - mean-r) ^ 2 ] real-price-history)

    ifelse (std-m > 0 and std-r > 0) [
      set price-correlation cov / (std-m * std-r)
    ] [
      set price-correlation 0
    ]

    ;; Finestra mobile sullo storico della correlazione
    set price-correlation-history lput price-correlation price-correlation-history
    if length price-correlation-history > window-size [
      set price-correlation-history but-first price-correlation-history
    ]

    set mean-price-correlation mean price-correlation-history
  ]

end

to aggiorna-plot-correlazione-prezzi

  carefully [
    set-current-plot "Correlazione Prezzi (Pearson r)"

    assicura-penna "zero" black
    plotxy day-index 0

    if btc-price > 0 [
      assicura-penna "correlazione" orange
      plotxy day-index price-correlation

      assicura-penna "correlazione-media" magenta
      plotxy day-index mean-price-correlation
    ]
  ] [ ]

end
\end{lstlisting}

La procedura \texttt{aggiorna-correlazione-prezzi} confronta l'andamento del
prezzo generato dal modello con quello del prezzo reale di Bitcoin. A ogni
aggiornamento, il prezzo endogeno \texttt{model-price} viene aggiunto allo storico

\[
\texttt{model-price-history},
\]

mentre il prezzo reale \texttt{btc-price} viene aggiunto allo storico

\[
\texttt{real-price-history}.
\]

Anche in questo caso viene usata una finestra mobile di ampiezza
\texttt{window-size}. Se gli storici superano questa lunghezza, il primo elemento
viene eliminato tramite \texttt{but-first}. In questo modo la correlazione misura
la somiglianza recente tra le due serie, invece di dipendere da tutta la storia
precedente della simulazione.

Quando sono disponibili almeno due osservazioni, il codice calcola la
correlazione di Pearson tra prezzi simulati e prezzi reali. Prima vengono
calcolate le medie:

\[
\bar{P}_m = \texttt{mean model-price-history},
\qquad
\bar{P}_r = \texttt{mean real-price-history}.
\]

Poi vengono calcolate la covarianza non normalizzata e le due deviazioni standard:

\[
\text{cov}
=
\sum_i
(P_{m,i}-\bar{P}_m)
(P_{r,i}-\bar{P}_r),
\]

\[
\sigma_m =
\sqrt{
\sum_i
(P_{m,i}-\bar{P}_m)^2
},
\qquad
\sigma_r =
\sqrt{
\sum_i
(P_{r,i}-\bar{P}_r)^2
}.
\]

La correlazione dei prezzi è quindi

\[
r_P =
\frac{\text{cov}}{\sigma_m\sigma_r}.
\]

Se una delle due deviazioni standard è nulla, la correlazione viene posta uguale
a \(0\). Questo evita divisioni per zero e indica che non è possibile valutare
una relazione lineare significativa se una delle due serie è costante.

Il valore ottenuto viene salvato in \texttt{price-correlation} e aggiunto allo
storico \texttt{price-correlation-history}. Anche questo storico viene mantenuto
entro la finestra mobile. La variabile \texttt{mean-price-correlation} contiene
la media recente delle correlazioni dei prezzi.

La procedura \texttt{aggiorna-plot-correlazione-prezzi} aggiorna il grafico
\texttt{Correlazione Prezzi (Pearson r)}. Il blocco \texttt{carefully} impedisce
che eventuali problemi con il plot interrompano la simulazione.

Nel grafico vengono tracciati tre elementi principali:

\[
\texttt{zero},
\qquad
\texttt{correlazione},
\qquad
\texttt{correlazione-media}.
\]

La penna \texttt{zero} disegna la linea di riferimento \(r=0\). Se il prezzo
reale è positivo, vengono poi disegnate la correlazione istantanea tra prezzi e
la sua media mobile. Il grafico permette quindi di osservare quanto il prezzo
endogeno del modello segua, almeno linearmente, il prezzo reale di Bitcoin.

\subsection{Plot di prezzo, volume e parametri d'ordine}

\begin{lstlisting}[caption={Aggiornamento dei plot e calcolo dei parametri d'ordine}]
;;;;;;;;;;;;;;;;;;;;;;;;;;;;;;;;;;;;;;
;; PLOT PREZZO E VOLUME
;;;;;;;;;;;;;;;;;;;;;;;;;;;;;;;;;;;;;;

to aggiorna-plot-prezzi-e-volume

  carefully [
    set-current-plot "Prezzo BTC (USD)"

    assicura-penna "prezzo-reale" blue
    plotxy day-index btc-price

    assicura-penna "prezzo-stimato" red
    plotxy day-index model-price

    set-current-plot "Volume BTC (USD)"

    assicura-penna "modello-usd" green
    plotxy day-index model-trading-volume-usd

    if real-trading-volume-usd > 0 [
      assicura-penna "reale-usd" gray
      plotxy day-index real-trading-volume-usd
    ]
  ] [ ]

end


;;;;;;;;;;;;;;;;;;;;;;;;;;;;;;;;;;;;;;
;; PARAMETRI D'ORDINE
;;;;;;;;;;;;;;;;;;;;;;;;;;;;;;;;;;;;;;

to-report order-parameter
  report order-parameter-of turtles
end

to-report order-parameter-a
  report order-parameter-of agents-a
end

to-report order-parameter-b
  report order-parameter-of agents-b
end

to-report order-parameter-c
  report order-parameter-of agents-c
end

to-report order-parameter-of [agenti-selezionati]

  if count agenti-selezionati = 0 [ report 0 ]

  let sx sum [dx] of agenti-selezionati
  let sy sum [dy] of agenti-selezionati

  report (sqrt (sx ^ 2 + sy ^ 2)) / count agenti-selezionati

end


;;;;;;;;;;;;;;;;;;;;;;;;;;;;;;;;;;;;;;
;; PLOT PARAMETRI D'ORDINE
;;;;;;;;;;;;;;;;;;;;;;;;;;;;;;;;;;;;;;

to aggiorna-plot-parametri

  carefully [
    set-current-plot "Parametri d'ordine"

    assicura-penna "globale" black
    plotxy ticks order-parameter

    assicura-penna "specie-a" yellow
    plotxy ticks order-parameter-a

    assicura-penna "specie-b" red
    plotxy ticks order-parameter-b

    assicura-penna "specie-c" cyan
    plotxy ticks order-parameter-c
  ] [ ]

end
\end{lstlisting}

La procedura \texttt{aggiorna-plot-prezzi-e-volume} aggiorna i grafici relativi
al prezzo e al volume. Il blocco è racchiuso in \texttt{carefully}, così eventuali
problemi con i plot o con le penne non interrompono la simulazione.

Nel plot

\[
\texttt{"Prezzo BTC (USD)"}
\]

vengono tracciate due serie: il prezzo reale di Bitcoin, \texttt{btc-price}, e il
prezzo stimato dal modello, \texttt{model-price}. La prima serie rappresenta il
dato esterno o storico, mentre la seconda rappresenta il prezzo generato dalla
dinamica degli agenti.

Nel plot

\[
\texttt{"Volume BTC (USD)"}
\]

viene invece tracciato il volume simulato in dollari, dato da
\texttt{model-trading-volume-usd}. Se il volume reale è disponibile e positivo,
viene aggiunta anche la serie \texttt{real-trading-volume-usd}. In questo modo è
possibile confrontare l'attività di mercato generata dal modello con quella
osservata nei dati reali.

La seconda parte del codice introduce i parametri d'ordine. Il parametro d'ordine
misura il grado di allineamento collettivo degli agenti. Il reporter generale

\[
\texttt{order-parameter}
\]

calcola il valore sull'intera popolazione, mentre i reporter
\texttt{order-parameter-a}, \texttt{order-parameter-b} e
\texttt{order-parameter-c} lo calcolano separatamente per le tre specie.

Il calcolo effettivo è svolto da

\[
\texttt{order-parameter-of [agenti-selezionati]}.
\]

Se l'insieme di agenti selezionato è vuoto, il reporter restituisce \(0\). In caso
contrario, vengono sommate le componenti direzionali degli agenti:

\[
S_x = \sum_i dx_i,
\qquad
S_y = \sum_i dy_i.
\]

Il parametro d'ordine è quindi

\[
\Phi =
\frac{\sqrt{S_x^2+S_y^2}}{N},
\]

dove \(N\) è il numero di agenti considerati. Se gli agenti sono orientati in
direzioni casuali, le componenti tendono a compensarsi e \(\Phi\) è vicino a
\(0\). Se invece gli agenti sono fortemente allineati, il modulo della somma
vettoriale cresce e \(\Phi\) si avvicina a \(1\).

Infine, la procedura \texttt{aggiorna-plot-parametri} aggiorna il grafico

\[
\texttt{"Parametri d'ordine"}.
\]

Nel plot vengono tracciati il parametro globale e i tre parametri specifici delle
specie A, B e C. Questo consente di osservare se l'allineamento emerge in tutta la
popolazione oppure solo all'interno di alcuni gruppi.

\subsection{Reporter di portafoglio aggregato}

\begin{lstlisting}[caption={Reporter per Bitcoin, liquidità e ricchezza complessiva degli agenti}]
;;;;;;;;;;;;;;;;;;;;;;;;;;;;;;;;;;;;;;
;; PORTAFOGLIO
;;;;;;;;;;;;;;;;;;;;;;;;;;;;;;;;;;;;;;

to-report total-bitcoin
  report sum [bitcoin-held] of turtles
end

to-report total-cash
  report sum [cash-balance] of turtles
end

to-report total-wealth
  let current-p model-price
  if current-p <= 0 [ report total-cash ]
  report total-cash + (total-bitcoin * current-p)
end

to-report bitcoin-a
  report sum [bitcoin-held] of agents-a
end

to-report bitcoin-b
  report sum [bitcoin-held] of agents-b
end

to-report bitcoin-c
  report sum [bitcoin-held] of agents-c
end
\end{lstlisting}

Questo blocco contiene alcuni reporter utili per misurare lo stato patrimoniale
complessivo degli agenti. Il reporter \texttt{total-bitcoin} calcola la quantità
totale di Bitcoin detenuta da tutti gli agenti:

\[
B_{\mathrm{tot}}=\sum_i B_i,
\]

dove \(B_i\) rappresenta la quantità di Bitcoin posseduta dall'agente \(i\).

Il reporter \texttt{total-cash} calcola invece la liquidità complessiva del
sistema:

\[
C_{\mathrm{tot}}=\sum_i C_i,
\]

dove \(C_i\) è il saldo monetario dell'agente \(i\).

Il reporter \texttt{total-wealth} restituisce la ricchezza totale degli agenti,
sommando la liquidità disponibile al valore in dollari dei Bitcoin posseduti:

\[
W_{\mathrm{tot}}
=
C_{\mathrm{tot}}
+
B_{\mathrm{tot}}P_m,
\]

dove \(P_m\) è il prezzo endogeno del modello. Nel codice, tale prezzo è
memorizzato nella variabile locale \texttt{current-p}. Se il prezzo non è
positivo, il reporter restituisce soltanto la liquidità totale, evitando di
valorizzare i Bitcoin con un prezzo nullo o negativo.

Gli ultimi tre reporter calcolano la quantità totale di Bitcoin detenuta
separatamente dalle tre specie:

\[
B_A = \sum_{i \in A} B_i,
\qquad
B_B = \sum_{i \in B} B_i,
\qquad
B_C = \sum_{i \in C} B_i.
\]

Nel codice questi valori corrispondono rispettivamente a \texttt{bitcoin-a},
\texttt{bitcoin-b} e \texttt{bitcoin-c}. Queste grandezze permettono di
confrontare la distribuzione dell'attività finanziaria tra i diversi gruppi di
agenti.

\subsection{Flocking dinamico con stubbornness variabile}

\begin{lstlisting}[caption={Procedura di flocking dinamico degli agenti}]
;;;;;;;;;;;;;;;;;;;;;;;;;;;;;;;;;;;;;;
;; FLOCKING DINAMICO
;;;;;;;;;;;;;;;;;;;;;;;;;;;;;;;;;;;;;;

to flock

  if any? flockmates [

    let eff-stubbornness dynamic-stubbornness

    let my-dx cos heading
    let my-dy sin heading

    let average-dx mean [cos heading] of flockmates
    let average-dy mean [sin heading] of flockmates

    let weighted-dx (eff-stubbornness * my-dx) + ((1 - eff-stubbornness) * average-dx)
    let weighted-dy (eff-stubbornness * my-dy) + ((1 - eff-stubbornness) * average-dy)

    if weighted-dx != 0 or weighted-dy != 0 [
      set heading atan weighted-dx weighted-dy
    ]

    let speed sqrt (velocity-x ^ 2 + velocity-y ^ 2)

    if speed > 0 [
      set velocity-x speed * cos heading
      set velocity-y speed * sin heading
    ]
  ]

end
\end{lstlisting}

La procedura \texttt{flock} aggiorna la direzione di ogni agente sulla base dei
suoi \texttt{flockmates}, cioè degli agenti che compongono la sua rete di
riferimento. Il blocco viene eseguito solo se l'agente possiede almeno un vicino:

\[
\texttt{if any? flockmates [ ... ]}.
\]

Il primo valore calcolato è

\[
\texttt{eff-stubbornness} = \texttt{dynamic-stubbornness}.
\]

Questa quantità misura quanto l'agente tende a mantenere la propria direzione
precedente. Una stubbornness alta rende l'agente più conservativo; una
stubbornness bassa lo rende più sensibile alla direzione media dei suoi
\texttt{flockmates}.

La direzione attuale dell'agente viene trasformata nelle sue componenti
vettoriali:

\[
d_x = \cos(\theta),
\qquad
d_y = \sin(\theta),
\]

dove \(\theta\) è il valore corrente di \texttt{heading}. Il codice calcola poi
anche la direzione media dei \texttt{flockmates}:

\[
\bar{d}_x = \left\langle \cos(\theta_j) \right\rangle,
\qquad
\bar{d}_y = \left\langle \sin(\theta_j) \right\rangle.
\]

Il nuovo orientamento è ottenuto come combinazione pesata tra la direzione
individuale dell'agente e la direzione media del gruppo:

\[
d_x^{\mathrm{new}}
=
s d_x
+
(1-s)\bar{d}_x,
\]

\[
d_y^{\mathrm{new}}
=
s d_y
+
(1-s)\bar{d}_y,
\]

dove \(s\) rappresenta la stubbornness efficace. Nel codice queste due componenti
sono \texttt{weighted-dx} e \texttt{weighted-dy}.

Se il vettore risultante non è nullo, il valore di \texttt{heading} viene
aggiornato tramite

\[
\texttt{set heading atan weighted-dx weighted-dy}.
\]

In questo modo l'agente non copia semplicemente la direzione media dei vicini,
ma si orienta verso una direzione intermedia tra il proprio stato precedente e lo
stato collettivo del gruppo.

Dopo aver aggiornato la direzione, la procedura calcola il modulo della velocità:

\[
v =
\sqrt{
\texttt{velocity-x}^2
+
\texttt{velocity-y}^2
}.
\]

Se la velocità è positiva, le componenti \texttt{velocity-x} e
\texttt{velocity-y} vengono ricostruite usando lo stesso modulo \(v\), ma con la
nuova direzione:

\[
v_x = v\cos(\theta),
\qquad
v_y = v\sin(\theta).
\]

Questa scelta modifica l'orientamento del moto senza alterare direttamente la
velocità scalare dell'agente. La procedura realizza quindi un meccanismo di
allineamento dinamico: gli agenti tendono a coordinarsi con la propria rete, ma
la forza dell'allineamento dipende dallo stato di stubbornness, che a sua volta
può variare con l'hype di mercato.

\subsection{Medie della posizione verticale delle specie}

\begin{lstlisting}[caption={Reporter per la posizione media lungo l'asse \(y\) delle specie A e B}]
;;;;;;;;;;;;;;;;;;;;;;;;;;;;;;;;;;;;;;
;; MEDIE DI POSIZIONE Y
;;;;;;;;;;;;;;;;;;;;;;;;;;;;;;;;;;;;;;

to-report mean-y-a
  if count agents-a = 0 [ report 0 ]
  report mean [ycor] of agents-a
end

to-report mean-y-b
  if count agents-b = 0 [ report 0 ]
  report mean [ycor] of agents-b
end
\end{lstlisting}

Questi due reporter calcolano la posizione media lungo l'asse verticale delle
specie \texttt{agents-a} e \texttt{agents-b}. La coordinata \(y\) è rilevante
perché, nel modello, distingue la zona associata all'acquisto da quella associata
alla vendita.

Il reporter \texttt{mean-y-a} restituisce

\[
\bar{y}_A =
\frac{1}{N_A}
\sum_{i \in A} y_i,
\]

dove \(N_A\) è il numero di agenti della specie A e \(y_i\) è la coordinata
verticale dell'agente \(i\). Prima del calcolo, il codice controlla che la specie
A non sia vuota:

\[
\texttt{if count agents-a = 0 [ report 0 ]}.
\]

Questo evita di calcolare una media su un insieme privo di agenti.

In modo analogo, il reporter \texttt{mean-y-b} calcola

\[
\bar{y}_B =
\frac{1}{N_B}
\sum_{i \in B} y_i.
\]

Queste quantità permettono di monitorare la posizione media delle due specie
nello spazio comportamentale. Se la media è positiva, la specie si trova
prevalentemente nella regione associata all'acquisto; se è negativa, si trova
prevalentemente nella regione associata alla vendita.

Il confronto tra \(\bar{y}_A\) e \(\bar{y}_B\) consente quindi di osservare se le
due popolazioni tendono a muoversi nella stessa direzione oppure se sviluppano
comportamenti opposti.

\subsection{Stima della pendenza tramite regressione lineare}

\begin{lstlisting}[caption={Reporter per il calcolo della pendenza di una regressione lineare}]
;;;;;;;;;;;;;;;;;;;;;;;;;;;;;;;;;;;;;;
;; REGRESSIONE LINEARE
;;;;;;;;;;;;;;;;;;;;;;;;;;;;;;;;;;;;;;

to-report linear-regression-slope [ ys ]

  let n length ys
  if n < 2 [ report 0 ]

  let xs n-values n [ i -> i ]
  let mean-x mean xs
  let mean-y mean ys

  let num sum (map [ [x y] -> (x - mean-x) * (y - mean-y) ] xs ys)
  let den sum (map [ x -> (x - mean-x) ^ 2 ] xs)

  if den = 0 [ report 0 ]

  report num / den

end
\end{lstlisting}

Il reporter \texttt{linear-regression-slope} calcola la pendenza della retta di
regressione lineare associata a una serie di valori \texttt{ys}. L'obiettivo è
stimare se una quantità osservata nel tempo tende mediamente ad aumentare,
diminuire oppure rimanere stabile.

La lista \texttt{ys} contiene i valori osservati, mentre le ascisse vengono
generate automaticamente tramite

\[
\texttt{let xs n-values n [ i -> i ]}.
\]

In questo modo, a una lista di \(n\) valori

\[
y_0,y_1,\dots,y_{n-1}
\]

viene associata la sequenza temporale

\[
x_0=0,\quad x_1=1,\quad \dots,\quad x_{n-1}=n-1.
\]

Se la lista contiene meno di due valori, il reporter restituisce \(0\), perché
non è possibile stimare una pendenza significativa con un solo punto.

Il codice calcola poi le medie delle ascisse e delle ordinate:

\[
\bar{x} = \frac{1}{n}\sum_i x_i,
\qquad
\bar{y} = \frac{1}{n}\sum_i y_i.
\]

La pendenza della regressione lineare viene calcolata come

\[
b =
\frac{
\sum_i (x_i-\bar{x})(y_i-\bar{y})
}{
\sum_i (x_i-\bar{x})^2
}.
\]

Nel codice, il numeratore è memorizzato nella variabile \texttt{num}, mentre il
denominatore è memorizzato in \texttt{den}. Se \texttt{den} è nullo, il reporter
restituisce \(0\), evitando una divisione non valida.

Il valore finale

\[
\texttt{report num / den}
\]

rappresenta quindi la pendenza stimata. Se la pendenza è positiva, la serie
\texttt{ys} mostra una tendenza crescente; se è negativa, mostra una tendenza
decrescente; se è vicina a zero, non presenta una tendenza lineare marcata.

\subsection{Correlazione di Pearson generica}

\begin{lstlisting}[caption={Reporter generico per il calcolo della correlazione di Pearson}]
;;;;;;;;;;;;;;;;;;;;;;;;;;;;;;;;;;;;;;
;; CORRELAZIONE DI PEARSON (GENERICA)
;;;;;;;;;;;;;;;;;;;;;;;;;;;;;;;;;;;;;;

to-report pearson-correlation [ xs ys ]

  let n length xs
  if n < 2 [ report 0 ]

  let mean-x mean xs
  let mean-y mean ys

  let cov sum (map [ [x y] -> (x - mean-x) * (y - mean-y) ] xs ys)
  let sx sqrt sum (map [ x -> (x - mean-x) ^ 2 ] xs)
  let sy sqrt sum (map [ y -> (y - mean-y) ^ 2 ] ys)

  if (sx = 0) or (sy = 0) [ report 0 ]

  report cov / (sx * sy)

end
\end{lstlisting}

Il reporter \texttt{pearson-correlation} calcola la correlazione di Pearson tra
due liste numeriche, indicate come \texttt{xs} e \texttt{ys}. Questo reporter è
generico, quindi può essere riutilizzato per confrontare serie temporali diverse,
come prezzi, volumi, parametri d'ordine o altre grandezze dinamiche del modello.

La variabile

\[
n = \texttt{length xs}
\]

rappresenta il numero di osservazioni considerate. Se la lista contiene meno di
due elementi, il reporter restituisce \(0\), perché non è possibile calcolare una
correlazione significativa con un solo dato.

Vengono poi calcolate le medie delle due serie:

\[
\bar{x} = \texttt{mean xs},
\qquad
\bar{y} = \texttt{mean ys}.
\]

La quantità \texttt{cov} rappresenta la covarianza non normalizzata tra le due
serie:

\[
\text{cov}
=
\sum_i
(x_i-\bar{x})(y_i-\bar{y}).
\]

Le quantità \texttt{sx} e \texttt{sy} rappresentano invece le deviazioni standard
non normalizzate:

\[
s_x =
\sqrt{
\sum_i (x_i-\bar{x})^2
},
\qquad
s_y =
\sqrt{
\sum_i (y_i-\bar{y})^2
}.
\]

Se una delle due serie è costante, allora una delle due deviazioni standard è
uguale a zero. In quel caso il reporter restituisce \(0\), evitando una divisione
per zero:

\[
\texttt{if (sx = 0) or (sy = 0) [ report 0 ]}.
\]

Se entrambe le serie hanno variazione non nulla, la correlazione di Pearson viene
calcolata come

\[
r =
\frac{\text{cov}}{s_x s_y}.
\]

Il valore restituito è compreso tra \(-1\) e \(1\). Un valore vicino a \(1\)
indica una relazione lineare positiva, un valore vicino a \(-1\) indica una
relazione lineare negativa, mentre un valore vicino a \(0\) indica assenza di una
relazione lineare evidente.

Nel modello, questa procedura è utile perché permette di misurare in modo
uniforme quanto due grandezze evolvano insieme nel tempo, senza dover riscrivere
ogni volta il calcolo della correlazione.

\subsection{Indice di eterogeneità della popolazione}

\begin{lstlisting}[caption={Reporter per l'indice di eterogeneità delle specie}]
;;;;;;;;;;;;;;;;;;;;;;;;;;;;;;;;;;;;;;
;; INDICE DI ETEROGENEITA'
;;;;;;;;;;;;;;;;;;;;;;;;;;;;;;;;;;;;;;

to-report heterogeneity-index

  let tot population-a + population-b + population-c
  if tot = 0 [ report 0 ]

  let sa population-a / tot
  let sb population-b / tot
  let sc population-c / tot

  report (sa ^ 2) + (sb ^ 2) + (sc ^ 2)

end
\end{lstlisting}

Il reporter \texttt{heterogeneity-index} calcola un indice basato sulla
distribuzione relativa delle tre popolazioni di agenti. La variabile

\[
\texttt{tot}
=
\texttt{population-a}
+
\texttt{population-b}
+
\texttt{population-c}
\]

rappresenta il numero complessivo di agenti previsto nel modello. Se questo
valore è nullo, il reporter restituisce \(0\), evitando divisioni non valide.

Successivamente vengono calcolate le quote relative delle tre specie:

\[
s_A =
\frac{\texttt{population-a}}{\texttt{tot}},
\qquad
s_B =
\frac{\texttt{population-b}}{\texttt{tot}},
\qquad
s_C =
\frac{\texttt{population-c}}{\texttt{tot}}.
\]

Il valore finale restituito è

\[
H =
s_A^2+s_B^2+s_C^2.
\]

Questo indice misura quanto la popolazione sia concentrata in poche specie.
Se le tre specie hanno dimensioni simili, le quote sono distribuite in modo più
uniforme e il valore di \(H\) è più basso. Se invece una specie domina sulle
altre, una delle quote è molto grande e il valore di \(H\) aumenta.

Nel caso di tre specie perfettamente bilanciate si ha

\[
s_A=s_B=s_C=\frac{1}{3},
\]

quindi

\[
H=
3\left(\frac{1}{3}\right)^2
=
\frac{1}{3}.
\]

Se invece tutta la popolazione appartiene a una sola specie, ad esempio

\[
s_A=1,
\qquad
s_B=s_C=0,
\]

allora

\[
H=1.
\]

Quindi il valore dell'indice varia tra \(1/3\) e \(1\): valori più bassi indicano
una popolazione più bilanciata tra le specie, mentre valori più alti indicano una
maggiore concentrazione in una singola specie.

\subsection{Aggiornamento delle metriche di stabilità}

\begin{lstlisting}[caption={Procedura di aggiornamento delle metriche di stabilità del sistema}]
;;;;;;;;;;;;;;;;;;;;;;;;;;;;;;;;;;;;;;
;; AGGIORNAMENTO METRICHE
;;;;;;;;;;;;;;;;;;;;;;;;;;;;;;;;;;;;;;

to aggiorna-metriche-stabilita

  set op-history lput order-parameter op-history
  if length op-history > window-size [ set op-history but-first op-history ]

  set ya-history lput mean-y-a ya-history
  if length ya-history > window-size [ set ya-history but-first ya-history ]

  set yb-history lput mean-y-b yb-history
  if length yb-history > window-size [ set yb-history but-first yb-history ]

  ifelse length op-history >= window-size [
    let log-op map [ v -> ln (v + 0.0001) ] op-history
    set lambda-instability linear-regression-slope log-op
  ][
    set lambda-instability 0
  ]

  ifelse length ya-history >= window-size [
    set phase-correlation-ab pearson-correlation ya-history yb-history
  ][
    set phase-correlation-ab 0
  ]

  set heterogeneity-index-value heterogeneity-index

end

to aggiorna-plot-stabilita
end
\end{lstlisting}

La procedura \texttt{aggiorna-metriche-stabilita} aggiorna alcune grandezze
utilizzate per descrivere la stabilità e la struttura collettiva del sistema.
Le tre serie storiche principali sono:

\[
\texttt{op-history},
\qquad
\texttt{ya-history},
\qquad
\texttt{yb-history}.
\]

La prima contiene lo storico del parametro d'ordine globale, mentre le altre due
contengono le posizioni medie lungo l'asse \(y\) delle specie A e B.

A ogni chiamata, il valore corrente del parametro d'ordine viene aggiunto a
\texttt{op-history}:

\[
\texttt{set op-history lput order-parameter op-history}.
\]

In modo analogo, \texttt{mean-y-a} e \texttt{mean-y-b} vengono aggiunti
rispettivamente a \texttt{ya-history} e \texttt{yb-history}. Per ciascuna lista
viene poi applicata una finestra mobile di ampiezza \texttt{window-size}: se la
lista diventa troppo lunga, il primo elemento viene eliminato con
\texttt{but-first}. In questo modo le metriche vengono calcolate solo sulla
dinamica recente del sistema.

Quando \texttt{op-history} ha raggiunto la lunghezza \texttt{window-size}, il
modello calcola una stima di instabilità:

\[
\lambda =
\text{slope}
\left[
\log(\Phi(t)+0.0001)
\right],
\]

dove \(\Phi(t)\) è il parametro d'ordine. Nel codice, la lista logaritmica è
costruita tramite

\[
\texttt{let log-op map [ v -> ln (v + 0.0001) ] op-history}.
\]

Il termine \(0.0001\) evita di calcolare il logaritmo di zero. La pendenza della
regressione lineare viene poi calcolata con

\[
\texttt{linear-regression-slope log-op}.
\]

Il risultato viene salvato in \texttt{lambda-instability}. Se questa quantità è
positiva, il parametro d'ordine tende a crescere nel tempo; se è negativa, tende
a diminuire; se è vicina a zero, il livello di ordine collettivo è
approssimativamente stabile.

Quando anche gli storici \texttt{ya-history} e \texttt{yb-history} hanno
raggiunto la lunghezza richiesta, il modello calcola la correlazione di Pearson
tra le posizioni medie verticali delle specie A e B:

\[
r_{AB}
=
\mathrm{corr}
\left(
\bar{y}_A,
\bar{y}_B
\right).
\]

Nel codice:

\[
\texttt{set phase-correlation-ab pearson-correlation ya-history yb-history}.
\]

Questa metrica misura se le due specie tendono a muoversi verticalmente in modo
coordinato. Un valore positivo indica movimenti mediamente concordi, un valore
negativo indica movimenti opposti, mentre un valore vicino a zero indica assenza
di una relazione lineare evidente.

Infine, il codice aggiorna l'indice di eterogeneità:

\[
\texttt{set heterogeneity-index-value heterogeneity-index}.
\]

Questa variabile conserva il valore dell'indice che misura quanto la popolazione
sia bilanciata o concentrata tra le tre specie.

La procedura \texttt{aggiorna-plot-stabilita} è invece attualmente vuota. Essa
può essere interpretata come un placeholder, cioè una procedura predisposta per
un eventuale futuro aggiornamento dei grafici relativi alle metriche di
stabilità.

% \begin{multicols}{2}
% % \noindent 
% Questo modello Agent-Based (ABM) simula le dinamiche di trading di un asset (ispirato al Bitcoin) utilizzando uno stormo di agenti divisi in tre categorie (Breed). Il modello evolve in due fasi: una \textbf{Fase 1 (Esogena)} guidata dai dati reali per calibrare il mercato, e una \textbf{Fase 2 (Endogena)} in cui il prezzo e i volumi sono determinati interamente dalle interazioni degli agenti.

% \section{Topologia e Comportamento degli Agenti}
% Gli agenti si muovono in uno spazio 2D che rappresenta il loro ``stato mentale'' rispetto al mercato:
% \begin{itemize}
%     \item \textbf{Asse Y (Sentiment):} Determina la direzione del trade. Se $Y \ge 0$, l'agente è propenso all'acquisto (Buy). Se $Y < 0$, è propenso alla vendita (Sell).
%     \item \textbf{Asse X (Convinzione):} Determina la probabilità effettiva di eseguire il trade. La probabilità $p$ cresce man mano che l'agente si sposta verso destra:
% \end{itemize}

% \begin{equation}
%     p = \frac{X - X_{min}}{X_{max} - X_{min}}
% \end{equation}

% \noindent Gli agenti sono divisi in tre tipologie di comportamento rispetto ai rendimenti (Return) del mercato:
% \begin{enumerate}
%     \item \textbf{Specie A (Gialli):} Trend-follower (seguono il mercato).
%     \item \textbf{Specie B (Rossi):} Contrarian (vanno contro il mercato).
%     \item \textbf{Specie C (Ciano):} Rumore di fondo (non influenzati direttamente dal rendimento per la direzione Y).
% \end{enumerate}

% \section{Dinamica del Movimento (Corrente e Hype)}
% Gli agenti vengono spinti da una ``corrente'' macroeconomica. La componente orizzontale (volatilità) spinge tutti verso una maggiore probabilità di trade:

% \begin{equation}
%     C_x = k \cdot \text{Volatilità}
% \end{equation}

% \noindent La componente verticale dipende dalla specie e dal rendimento (Return):
% \begin{align}
%     C_y^{(A)} &= h \cdot \text{Return} \\
%     C_y^{(B)} &= -h \cdot \text{Return} \\
%     C_y^{(C)} &= 0
% \end{align}

% \subsection*{Hype e Stubbornness (Testardaggine)}
% L'Hype (stress di mercato) è definito come la somma del rendimento assoluto e della volatilità:
% \begin{equation}
%     H = |\text{Return}| + \text{Volatilità}
% \end{equation}

% \noindent Quando l'Hype sale, gli agenti diventano meno ``testardi'' ($S_{dyn}$) e più inclini a seguire il gregge (flocking):
% \begin{equation}
%     S_{dyn} = \max\left(S_{min}, \min\left(S_{base}, S_{base} - (K_{hype} \cdot H)\right)\right)
% \end{equation}
% Dove $K_{hype}$ è la sensibilità all'hype impostata dall'utente.

% \section{Transizione alla Fase 2 e Prezzo Endogeno}
% Il modello calcola la correlazione di Pearson $r$ tra il volume generato dagli agenti e il volume reale. Se $r \ge \text{Threshold}$ per più di 5 tick consecutivi, il modello ``sgancia'' i dati reali e inizia a generare il prezzo endogenamente.

% \vspace{0.3cm}
% \noindent Il nuovo prezzo viene calcolato in base alla \textbf{Pressione di Mercato} ($P_{m}$), definita come la media delle coordinate Y degli agenti rispetto al massimo possibile:
% \begin{equation}
%     P_{m} = \frac{\bar{Y}}{Y_{max}}
% \end{equation}

% \noindent Il prezzo endogeno al tempo $t+1$ si aggiorna usando il parametro \texttt{price-sensitivity} ($S_p$):
% \begin{equation}
%     \text{Price}_{t+1} = \max\left(1, \text{Price}_t \cdot (1 + S_p \cdot P_{m})\right)
% \end{equation}

% \section{Il Doppio Motore di Isteresi (Memoria di Mercato)}
% L'isteresi simula il fatto che i volumi in dollari non dipendono solo dall'azione istantanea, ma mantengono ``memoria'' dei picchi storici raggiunti (All-Time High). Nel modello agiscono due forze di isteresi moltiplicative.

% \subsection{Isteresi del Prezzo ($\beta_{vol}$)}
% Misura quanto il mercato ricorda i volumi basandosi sul crollo del prezzo rispetto al suo picco. Siano:
% \begin{equation}
%     R_{p}^{(current)} = \frac{\text{Price}_{corrente}}{\text{Price}_{iniziale}}, \quad R_{p}^{(ath)} = \frac{\text{Price}_{ath}}{\text{Price}_{iniziale}}
% \end{equation}
% Il fattore di isteresi del prezzo $F_p$ è:
% \begin{equation}
%     F_p = \max \left( 0.01, R_{p}^{(current)} + \beta_{vol} \left( R_{p}^{(ath)} - R_{p}^{(current)} \right) \right)
% \end{equation}

% \subsection{Isteresi del Volume ($\gamma_{vol}$)}
% Misura quanto il mercato mantiene i volumi basandosi puramente sulle masse scambiate in passato (memoria dei flussi). Siano:
% \begin{equation}
%     R_{v}^{(current)} = \frac{\text{Volume}_{corrente}}{\text{Volume}_{iniziale}}, \quad R_{v}^{(ath)} = \frac{\text{Volume}_{ath}}{\text{Volume}_{iniziale}}
% \end{equation}
% Il fattore di isteresi del volume $F_v$ è:
% \begin{equation}
%     F_v = \max \left( 0.01, R_{v}^{(current)} + \gamma_{vol} \left( R_{v}^{(ath)} - R_{v}^{(current)} \right) \right)
% \end{equation}

% \subsection{Volume Finale in Dollari}
% Il volume totale scambiato in USD dal modello, che tiene conto del volume grezzo scambiato dagli agenti ($Vol_{raw}$) moltiplicato per il prezzo attuale e ``gonfiato'' dalla memoria del mercato, è infine:
% \begin{equation}
%     \text{Vol}_{USD} = \text{Vol}_{raw} \cdot \text{Price}_{corrente} \cdot F_p \cdot F_v
% \end{equation}

\noindent
Il progetto sviluppa un modello \textit{agent-based} per simulare la formazione
di dinamiche di trading su un asset digitale ispirato a Bitcoin. L'obiettivo non
è costruire un modello predittivo del prezzo reale, ma realizzare un ambiente
computazionale in cui il comportamento collettivo degli agenti possa generare
grandezze aggregate confrontabili con quelle osservate nei dati di mercato:
prezzo, volume, volatilità, correlazioni e grado di coordinamento della
popolazione.

Il modello combina due livelli descrittivi. Da un lato è presente una componente
finanziaria, costituita da prezzo, rendimento, volatilità, volume di scambio e
portafoglio degli agenti. Dall'altro lato è presente una componente dinamica di
tipo \textit{flocking}, in cui gli agenti si muovono in uno spazio bidimensionale,
interagiscono con i propri vicini e modificano il proprio comportamento in
funzione dello stato del mercato.

Una caratteristica centrale del modello è la presenza di due fasi. Nella
\textbf{Fase 1}, detta esogena, il prezzo del modello viene agganciato al prezzo
reale o sintetico fornito dalla serie storica. In questa fase il sistema usa il
dato esterno come riferimento iniziale. Nella \textbf{Fase 2}, detta endogena, il
prezzo non viene più copiato dai dati, ma viene prodotto dalla dinamica interna
degli agenti. Il passaggio tra le due fasi avviene automaticamente quando il
volume generato dal modello risulta sufficientemente correlato con il volume
reale.

\section{Struttura generale del modello}

Il modello considera tre popolazioni di agenti, indicate come \texttt{agents-a},
\texttt{agents-b} e \texttt{agents-c}. Le tre specie non rappresentano singoli
investitori reali, ma categorie comportamentali astratte. Esse permettono di
studiare come gruppi con risposte diverse allo stesso segnale di mercato possano
produrre dinamiche aggregate differenti.

Ogni agente possiede alcune variabili individuali fondamentali:
\begin{itemize}
    \item una posizione nello spazio bidimensionale, data dalle coordinate
    \texttt{xcor} e \texttt{ycor};
    \item una velocità, descritta dalle componenti \texttt{velocity-x} e
    \texttt{velocity-y};
    \item un insieme di vicini o \texttt{flockmates}, che definisce la rete di
    riferimento per il meccanismo di allineamento;
    \item un portafoglio finanziario, composto da una quantità di Bitcoin
    detenuta, \texttt{bitcoin-held}, e da una disponibilità liquida,
    \texttt{cash-balance}.
\end{itemize}

Durante il \texttt{setup}, gli agenti vengono creati in una regione circolare
centrale, ricevono una direzione iniziale casuale e vengono inizializzati con una
dotazione casuale di Bitcoin e liquidità. Inoltre, per ogni specie viene costruita
una rete fissa di interazione: ogni agente seleziona un certo numero di
\texttt{flockmates} appartenenti alla stessa specie. In questo modo il modello
separa la vicinanza sociale, rappresentata dalla rete, dalla vicinanza spaziale,
rappresentata dalla posizione nello spazio.

\section{Interpretazione dello spazio comportamentale}

Lo spazio bidimensionale in cui si muovono gli agenti non rappresenta uno spazio
fisico, ma uno spazio comportamentale. Le coordinate dell'agente sintetizzano il
suo stato rispetto al mercato.

L'asse verticale \(y\) rappresenta la direzione dell'azione di trading:
\begin{itemize}
    \item se \(y \geq 0\), l'agente si trova nella regione di acquisto;
    \item se \(y < 0\), l'agente si trova nella regione di vendita.
\end{itemize}

L'asse orizzontale \(x\) rappresenta invece la propensione operativa dell'agente,
cioè la probabilità che esso esegua effettivamente una compravendita. Tale
probabilità è calcolata normalizzando la posizione orizzontale dell'agente
rispetto alla larghezza del mondo:

\begin{equation}
    p_i =
    \frac{x_i - x_{\min}}{x_{\max} - x_{\min}}.
\end{equation}

Nel codice il valore viene ulteriormente limitato all'intervallo \([0,1]\), in
modo da ottenere sempre una probabilità valida:

\begin{equation}
    p_i =
    \min\left[
    1,
    \max\left(
    0,
    \frac{x_i - x_{\min}}{x_{\max} - x_{\min}}
    \right)
    \right].
\end{equation}

Questa scelta permette di attribuire un significato comportamentale semplice alle
coordinate: la posizione verticale stabilisce se l'agente tende a comprare o a
vendere, mentre la posizione orizzontale stabilisce con quale probabilità
partecipa al mercato.

\section{Specie di agenti e risposta al rendimento}

Le tre specie di agenti differiscono per il modo in cui rispondono al rendimento
del mercato. Il rendimento attivo può essere quello reale, nella Fase 1, oppure
quello endogeno, nella Fase 2. Indichiamo tale rendimento con \(r(t)\).

La componente verticale della corrente è definita in modo diverso per le tre
specie:

\begin{align}
    J_y^{(A)}(t) &= h r(t), \\
    J_y^{(B)}(t) &= -h r(t), \\
    J_y^{(C)}(t) &= 0.
\end{align}

La specie A può essere interpretata come una popolazione \textit{trend-following}:
quando il rendimento è positivo, essa viene spinta verso la regione di acquisto;
quando il rendimento è negativo, viene spinta verso la regione di vendita.

La specie B ha una risposta opposta e può essere interpretata come una popolazione
\textit{contrarian}: tende a muoversi contro il segnale di rendimento. La specie C
non riceve una spinta verticale diretta dal rendimento e funziona come componente
neutrale o rumore di fondo.

Questa distinzione consente di osservare come l'interazione tra gruppi con
comportamenti diversi influenzi la dinamica complessiva del mercato simulato.

\section{Corrente di mercato e dinamica del movimento}

Il movimento degli agenti è influenzato da una corrente di mercato. La corrente
ha una componente orizzontale e una componente verticale. La componente
orizzontale dipende dalla volatilità attiva:

\begin{equation}
    J_x(t) = k \sigma(t),
\end{equation}

dove \(k\) è un parametro di intensità e \(\sigma(t)\) è la volatilità. Nella
Fase 1 viene usata la volatilità reale, mentre nella Fase 2 viene usata la
volatilità endogena del modello.

La funzione della componente orizzontale è spingere gli agenti lungo l'asse della
propensione operativa. Quando la volatilità è più alta, il mercato è più
instabile e gli agenti tendono a muoversi verso regioni in cui la probabilità di
trading è maggiore.

Il movimento è descritto mediante una velocità vettoriale. A ogni passo, la
corrente modifica le componenti della velocità:

\begin{align}
    v_x(t+1) &= v_x(t) + J_x(t), \\
    v_y(t+1) &= v_y(t) + J_y(t).
\end{align}

La nuova posizione prevista è poi calcolata come

\begin{align}
    x_i(t+1) &= x_i(t) + v_{x,i}(t+1), \\
    y_i(t+1) &= y_i(t) + v_{y,i}(t+1).
\end{align}

Il modello utilizza condizioni al bordo riflettenti. Se un agente raggiunge un
bordo del mondo, la componente corrispondente della velocità viene invertita.
Questa scelta mantiene gli agenti confinati nello spazio comportamentale ed evita
che escano dal dominio della simulazione.

Infine, viene introdotta una perturbazione casuale della direzione di moto. La
velocità viene ruotata di un angolo casuale compreso tra
\(-\texttt{temperature}\) e \(+\texttt{temperature}\). Il parametro
\texttt{temperature} controlla quindi il grado di rumore della dinamica: valori
bassi producono traiettorie più regolari, mentre valori alti rendono il moto più
disordinato.

\section{Flocking, hype e stubbornness}

\subsection{Condizioni iniziali}
La scelta di inizializzare gli agenti in una regione circolare centrale non deve
essere interpretata come una descrizione realistica dello stato iniziale del
mercato, ma come una condizione iniziale neutra. Essa rappresenta una popolazione
inizialmente poco polarizzata, nella quale gli agenti non sono ancora fortemente
orientati né verso l'acquisto né verso la vendita.

Questa scelta è utile per studiare se le dinamiche collettive osservate nel
modello emergano dalle regole di interazione, dalla corrente di mercato e dal
meccanismo di flocking, piuttosto che essere imposte artificialmente dalla
condizione iniziale.

Tuttavia, la condizione iniziale non è unica. Per analizzare la robustezza del
modello è opportuno confrontare diverse inizializzazioni: distribuzione centrale,
distribuzione uniforme, distribuzioni sbilanciate verso l'acquisto o la vendita,
e condizioni differenziate per specie. In questo modo è possibile distinguere gli
effetti realmente generati dalla dinamica del modello da quelli dipendenti dalla
scelta iniziale.\\


Gli agenti non si muovono in modo completamente indipendente. Ogni agente tende
ad allineare la propria direzione con quella dei propri \texttt{flockmates}. Il
nuovo orientamento viene calcolato come combinazione pesata tra la direzione
individuale dell'agente e la direzione media dei suoi vicini.

Sia \(\theta_i\) la direzione dell'agente e siano

\begin{equation}
    d_{x,i} = \cos(\theta_i),
    \qquad
    d_{y,i} = \sin(\theta_i)
\end{equation}

le componenti della sua direzione. Indicando con
\(\bar{d}_{x,i}\) e \(\bar{d}_{y,i}\) le componenti medie dei suoi
\texttt{flockmates}, la direzione aggiornata è costruita come

\begin{align}
    d_{x,i}^{new}
    &=
    s(t)d_{x,i}
    +
    \left[1-s(t)\right]\bar{d}_{x,i}, \\
    d_{y,i}^{new}
    &=
    s(t)d_{y,i}
    +
    \left[1-s(t)\right]\bar{d}_{y,i}.
\end{align}

Il parametro \(s(t)\) rappresenta la \textit{stubbornness}, cioè la tendenza
dell'agente a mantenere la propria direzione precedente. Se \(s(t)\) è alto,
l'agente è poco influenzato dal gruppo; se \(s(t)\) è basso, l'agente tende ad
allinearsi maggiormente ai propri vicini.

Nel modello la stubbornness non è costante, ma dipende dall'hype di mercato. Lo
hype è definito come

\begin{equation}
    H(t)=|r_m(t)|+\sigma_m(t),
\end{equation}

dove \(r_m(t)\) è il rendimento endogeno del modello e \(\sigma_m(t)\) è la
volatilità endogena. Quando il mercato è più instabile, cioè quando rendimento e
volatilità sono più elevati, gli agenti diventano meno testardi e più sensibili
al comportamento collettivo.

La stubbornness dinamica è calcolata come

\begin{equation}
    s(t)
    =
    \max
    \left[
    s_{\min},
    \min
    \left(
    s_{\mathrm{base}},
    s_{\mathrm{base}} -
    K_{\mathrm{hype}}H(t)
    \right)
    \right].
\end{equation}

Il parametro \(K_{\mathrm{hype}}\) controlla quanto rapidamente la stubbornness
diminuisce all'aumentare dello stress di mercato.

\section{Trading e portafoglio degli agenti}

Ogni agente possiede un portafoglio formato da liquidità e Bitcoin. La procedura
di trading utilizza il prezzo endogeno corrente del modello, indicato con
\(P_m(t)\). Se il prezzo è positivo, l'agente può comprare o vendere a seconda
della propria posizione verticale.

Se \(y_i \geq 0\), l'agente si trova nella regione di acquisto. La quantità
massima acquistabile è

\begin{equation}
    q_i^{buy}
    =
    \min
    \left(
    \texttt{trade-size},
    \left\lfloor
    \frac{C_i}{P_m(t)}
    \right\rfloor
    \right),
\end{equation}

dove \(C_i\) è la liquidità disponibile dell'agente. Questa formula impedisce
all'agente di comprare più Bitcoin di quanto possa permettersi.

Se \(y_i < 0\), l'agente si trova nella regione di vendita. La quantità massima
vendibile è

\begin{equation}
    q_i^{sell}
    =
    \min
    \left(
    \texttt{trade-size},
    B_i
    \right),
\end{equation}

dove \(B_i\) è la quantità di Bitcoin detenuta dall'agente. In questo modo non è
possibile vendere Bitcoin non posseduti.

L'operazione avviene solo se un numero casuale uniforme tra 0 e 1 è minore della
probabilità di scambio \(p_i\). Ogni acquisto o vendita riuscita contribuisce al
volume di trading del modello, misurato inizialmente in Bitcoin.

\section{Dati di mercato e scenari}

Il modello può essere eseguito utilizzando diverse sorgenti di dati. La procedura
\texttt{scarica-dati} seleziona lo scenario in base alla variabile
\texttt{scenario}. Le possibilità previste sono:

\begin{itemize}
    \item caricamento di una serie storica da file CSV;
    \item utilizzo di un prezzo live tramite API;
    \item generazione di scenari sintetici di mercato.
\end{itemize}

Gli scenari sintetici includono:
\begin{itemize}
    \item \textit{bull market}, con drift positivo;
    \item \textit{bear market}, con drift negativo;
    \item mercato laterale, con drift nullo;
    \item crash improvviso;
    \item bolla speculativa seguita da collasso.
\end{itemize}

Questa struttura rende il modello flessibile: è possibile confrontare la risposta
degli agenti sia a dati esterni sia a condizioni artificiali controllate.

\section{Fase 1: dinamica esogena}

Nella Fase 1 il modello è guidato dai dati esterni. Il prezzo endogeno viene
posto uguale al prezzo reale o sintetico letto dalla serie storica:

\begin{equation}
    P_m(t) = P_{BTC}(t).
\end{equation}

Anche rendimento e volatilità del modello vengono copiati dai dati esterni:

\begin{align}
    r_m(t) &= r_{BTC}(t), \\
    \sigma_m(t) &= \sigma_{BTC}(t).
\end{align}

Questa fase ha la funzione di inizializzare il sistema e permettere agli agenti
di produrre volumi confrontabili con quelli della serie di riferimento. Il
modello resta in questa fase finché la correlazione media tra volume simulato e
volume reale non supera stabilmente una soglia prefissata.

\section{Correlazione dei volumi e passaggio alla Fase 2}

Il modello confronta il volume reale con il volume generato dagli agenti mediante
la correlazione di Pearson. Date due serie \(X\) e \(Y\), la correlazione è

\begin{equation}
    r =
    \frac{
    \sum_i (x_i-\bar{x})(y_i-\bar{y})
    }{
    \sqrt{\sum_i (x_i-\bar{x})^2}
    \sqrt{\sum_i (y_i-\bar{y})^2}
    }.
\end{equation}

Nel caso dei volumi, \(X\) rappresenta lo storico recente dei volumi simulati,
mentre \(Y\) rappresenta lo storico recente dei volumi reali. Il calcolo viene
effettuato su una finestra mobile di ampiezza \texttt{window-size}, così da
misurare la somiglianza recente tra le due serie.

Il passaggio alla Fase 2 avviene quando la correlazione media dei volumi supera
la soglia \texttt{pearson-threshold} per più di cinque tick consecutivi:

\begin{equation}
    \overline{r}_V \geq \texttt{pearson-threshold}.
\end{equation}

Quando questa condizione viene soddisfatta, la variabile
\texttt{current-disabled?} viene impostata a \texttt{true}. Da quel momento il
modello smette di copiare il prezzo esterno e inizia a generare il prezzo in modo
endogeno.

\section{Fase 2: prezzo endogeno}

Nella Fase 2 il prezzo viene determinato dalla posizione media verticale degli
agenti. La coordinata \(y\) è interpretata come indicatore della pressione di
acquisto o vendita. La pressione di mercato è definita come

\begin{equation}
    M(t)=\frac{\bar{y}(t)}{y_{\max}},
\end{equation}

dove \(\bar{y}(t)\) è la media delle coordinate verticali di tutti gli agenti e
\(y_{\max}\) è il massimo valore verticale del mondo NetLogo.

Se \(M(t)>0\), la popolazione è mediamente orientata verso l'acquisto e il prezzo
subisce una pressione rialzista. Se \(M(t)<0\), la popolazione è mediamente
orientata verso la vendita e il prezzo subisce una pressione ribassista.

Il prezzo endogeno viene aggiornato secondo

\begin{equation}
    P_m(t+1)
    =
    \max
    \left[
    1,
    P_m(t)
    \left(
    1+
    S_p M(t)
    \right)
    \right],
\end{equation}

dove \(S_p=\texttt{price-sensitivity}\) controlla la sensibilità del prezzo alla
pressione di mercato. Il limite inferiore pari a 1 evita che il prezzo diventi
nullo o negativo.

Il rendimento endogeno viene poi calcolato come

\begin{equation}
    r_m(t)
    =
    \frac{P_m(t)-P_m(t-1)}{P_m(t-1)}.
\end{equation}

La volatilità endogena è approssimata dal valore assoluto del rendimento, con una
soglia minima:

\begin{equation}
    \sigma_m(t)
    =
    \max
    \left(
    0.005,
    |r_m(t)|
    \right).
\end{equation}

La soglia minima impedisce alla volatilità di annullarsi completamente e mantiene
attivi i meccanismi dipendenti da essa.

\section{Volume simulato e isteresi}

Il volume grezzo del modello, \(\text{Vol}_{raw}\), è dato dalla quantità totale
di Bitcoin scambiata dagli agenti nel tick corrente. Per confrontarlo con un
volume reale espresso in dollari, esso viene moltiplicato per il prezzo endogeno:

\begin{equation}
    \text{Vol}_{raw}^{USD}
    =
    \text{Vol}_{raw} \cdot P_m(t).
\end{equation}

Nel modello viene introdotto un fattore di memoria aggregata del mercato. Tale
fattore è ispirato al concetto di isteresi, intesa come dipendenza dello stato
attuale di un sistema non solo dal valore corrente delle variabili rilevanti, ma
anche dal percorso con cui tale stato è stato raggiunto.

Nel caso considerato, il volume simulato non dipende soltanto dal prezzo corrente,
ma anche dal massimo prezzo raggiunto nella finestra di simulazione. Due stati con
lo stesso prezzo corrente possono quindi produrre livelli diversi di attività di
mercato se sono stati preceduti da storie differenti. In questo senso il termine
può essere interpretato come una forma fenomenologica di path dependence.

È importante sottolineare che tale memoria non rappresenta necessariamente una
memoria esplicita dei singoli agenti. Nel modello attuale essa è introdotta come
variabile macroscopica aggregata, utile a rappresentare la persistenza
dell'attenzione e dell'attività di mercato dopo fasi di forte crescita del prezzo.
Siano

\begin{equation}
    R_{\mathrm{current}}
    =
    \frac{P_m(t)}{P_{\mathrm{initial}}},
    \qquad
    R_{\mathrm{ATH}}
    =
    \frac{P_{\mathrm{ATH}}}{P_{\mathrm{initial}}},
\end{equation}

dove \(P_{\mathrm{initial}}\) è il prezzo iniziale e \(P_{\mathrm{ATH}}\) è il
massimo prezzo raggiunto nella finestra di simulazione. Il fattore di isteresi è

\begin{equation}
    H_V(t)
    =
    \max
    \left[
    0.01,
    R_{\mathrm{current}}
    +
    \beta_{\mathrm{vol}}
    \left(
    R_{\mathrm{ATH}} -
    R_{\mathrm{current}}
    \right)
    \right].
\end{equation}

Il parametro \(\beta_{\mathrm{vol}}\) controlla il peso della memoria del massimo
storico. Se \(\beta_{\mathrm{vol}}\) è vicino a zero, il fattore dipende quasi
solo dal prezzo corrente. Se invece è più alto, il modello mantiene una memoria
più marcata dei picchi passati.

Il volume finale in dollari prodotto dal modello è quindi

\begin{equation}
    \text{Vol}_{model}^{USD}
    =
    \text{Vol}_{raw}
    \cdot
    P_m(t)
    \cdot
    H_V(t).
\end{equation}

\section{Metriche collettive e stabilità}

Oltre a prezzo e volume, il modello calcola alcune metriche utili per analizzare
la struttura collettiva degli agenti.

La prima è il parametro d'ordine, che misura il grado di allineamento della
popolazione. Per un insieme di \(N\) agenti, esso è definito come

\begin{equation}
    \Phi
    =
    \frac{
    \sqrt{
    \left(\sum_i dx_i\right)^2
    +
    \left(\sum_i dy_i\right)^2
    }
    }{N}.
\end{equation}

Se gli agenti sono orientati in direzioni casuali, le componenti vettoriali si
compensano e \(\Phi\) tende a valori bassi. Se invece gli agenti sono allineati,
\(\Phi\) si avvicina a 1. Il modello calcola sia il parametro d'ordine globale,
sia quello specifico per ciascuna specie.

Una seconda metrica è

\begin{equation}
    \lambda =
    \text{slope}
    \left[
    \log(\Phi(t)+0.0001)
    \right],
\end{equation}

dove la pendenza è calcolata tramite regressione lineare su una finestra mobile.
Questa grandezza fornisce una misura della tendenza recente del parametro
d'ordine: se \(\lambda>0\), l'ordine collettivo tende ad aumentare; se
\(\lambda<0\), tende a diminuire.

Il modello calcola inoltre la correlazione di fase tra le specie A e B,
utilizzando le loro posizioni medie verticali:

\begin{equation}
    r_{AB}
    =
    \mathrm{corr}
    \left(
    \bar{y}_A,
    \bar{y}_B
    \right).
\end{equation}

Questa quantità misura se le due specie si muovono nella stessa direzione lungo
l'asse comportamentale di acquisto-vendita oppure se tendono a muoversi in modo
opposto.

Infine, viene calcolato un indice basato sulla distribuzione relativa delle tre
popolazioni:

\begin{equation}
    H =
    s_A^2+s_B^2+s_C^2,
\end{equation}

dove \(s_A\), \(s_B\) e \(s_C\) sono le quote relative delle tre specie. Questo
indice è minimo quando le tre popolazioni sono bilanciate e aumenta quando una
specie domina sulle altre. Per questo motivo può essere interpretato come un
indice di concentrazione della composizione della popolazione.

\section{Sintesi del funzionamento del modello}

Il funzionamento complessivo del modello può essere riassunto come segue. In una
prima fase, gli agenti vengono inizializzati con portafogli eterogenei e si
muovono in uno spazio comportamentale influenzato da segnali di mercato esterni.
La loro posizione determina la probabilità e la direzione delle operazioni di
trading. Le operazioni aggregate producono un volume simulato, che viene
confrontato con il volume reale tramite correlazione di Pearson.

Quando la correlazione dei volumi supera stabilmente la soglia impostata, il
modello passa alla fase endogena. Da quel momento il prezzo non è più imposto
dall'esterno, ma emerge dalla distribuzione degli agenti nello spazio. Il mercato
simulato diventa quindi il risultato dell'interazione tra movimento, flocking,
portafogli, segnali di rendimento, volatilità e memoria dei picchi passati.

Il modello costituisce quindi un laboratorio computazionale per studiare come
semplici regole individuali possano generare dinamiche aggregate complesse e
come un sistema inizialmente guidato dai dati possa evolvere verso una dinamica
autonoma.
% \end{multicols}

% =========================================================
% BIBLIOGRAFIA
% =========================================================
% Scommenta questa riga se vuoi stampare la bibliografia finale.
% \printbibliography

\end{document}
